\documentclass[12pt,openany]{ctexbook}

% 上海财经大学硕士学位论文版式
% 版式目标：A4，黑白色正文，左侧装订，三线表，表题在上，图题在下
% 生效源文件：chapters_new.tex（第1-5章）、chapters_6_9.tex（第6-7章与附录）、backmatter.tex。
% chapter7_new.tex、chapter89_new.tex、appendix_new.tex 不参与编译，仅作历史存档。
\usepackage[a4paper,top=3.8cm,bottom=3.8cm,left=3.2cm,right=3.2cm,headheight=13pt,headsep=0.35cm,footskip=0.8cm]{geometry}
\usepackage{setspace}
\usepackage{amsmath,amssymb,amsthm}
\usepackage{booktabs}
\usepackage[table]{xcolor}
\usepackage{tcolorbox}
\tcbuselibrary{breakable}
\usepackage{graphicx}
\usepackage{enumitem}
\usepackage{array}
\usepackage{caption}
\usepackage{tocloft}
\renewcommand{\cftchapfont}{\mdseries\songti\zihao{4}}
\renewcommand{\cftchappagefont}{\mdseries\songti\zihao{4}}
\renewcommand{\cftsecfont}{\mdseries\songti\zihao{-4}}
\renewcommand{\cftsecpagefont}{\mdseries\songti\zihao{-4}}
\renewcommand{\cftsubsecfont}{\mdseries\songti\zihao{5}}
\renewcommand{\cftsubsecpagefont}{\mdseries\songti\zihao{5}}
\setlength{\cftchapindent}{2em}
\setlength{\cftsecindent}{4em}
\setlength{\cftsecnumwidth}{4.2em}
\renewcommand{\cftsecaftersnum}{\hspace{1em}}
\setlength{\cftsubsecindent}{7em}
\setlength{\cftsubsecnumwidth}{2.5em}
\renewcommand{\cftsubsecaftersnum}{\hspace{0.5em}}
\usepackage{fancyhdr}
\usepackage[hidelinks]{hyperref}
\usepackage[round,authoryear]{natbib}
\bibpunct{(}{)}{;}{a}{,}{,}
\renewcommand{\bibname}{参考文献}
\usepackage{fontspec}
\usepackage{seqsplit}
\usepackage{tikz}
\usepackage{algorithm}
\usepackage{float}
\usepackage{placeins}
\usepackage{algpseudocode}
\usepackage{listings}
\lstset{
  basicstyle=\ttfamily\footnotesize,
  breaklines=true,
  columns=fullflexible,
  frame=single,
  numbers=left,
  numberstyle=\tiny,
  showstringspaces=false,
  keepspaces=true,
  tabsize=2
}
\usetikzlibrary{arrows.meta,positioning,fit,backgrounds}

\newfontfamily\arial{Arial}
\setmainfont{Times New Roman}
\setCJKmainfont[AutoFakeBold=2,BoldFont=SimHei,ItalicFont=SimSun]{SimSun}
\setCJKsansfont[AutoFakeBold=2,BoldFont=SimHei]{SimHei}
\setCJKmonofont{SimSun}
\setCJKfamilyfont{song}[AutoFakeBold=2,BoldFont=SimHei]{SimSun}
\setCJKfamilyfont{hei}[AutoFakeBold=2,BoldFont=SimHei]{SimHei}
\setCJKfamilyfont{zhsong}[AutoFakeBold=2,BoldFont=SimHei]{SimSun}
\setCJKfamilyfont{zhhei}[AutoFakeBold=2,BoldFont=SimHei]{SimHei}
\setCJKfamilyfont{zhfs}[AutoFakeBold=2,BoldFont=SimHei]{FangSong}
\setCJKfamilyfont{callig}{STXingkai}
\newcommand{\calligfont}{\CJKfamily{callig}}
\XeTeXgenerateactualtext=1

% 全文保持黑色，不使用彩色链接或彩色表格
\hypersetup{colorlinks=false,hidelinks,linkcolor=black,citecolor=black,urlcolor=black,
  pdftitle={FinPlanRAG：面向金融混合查询的契约路由与锚点演化框架}}
% 预答辩/查重匿名版：不写入 pdfauthor，终稿恢复 pdfauthor={高深深}

% 表题在上，图题在下；图序/表序与题名之间空两个字符
\DeclareCaptionFont{sufeten}{\fontsize{10.5pt}{12.6pt}\selectfont}
\DeclareCaptionLabelSeparator{sufetwo}{\hspace{2em}}
\captionsetup[table]{position=top,font=sufeten,labelfont=bf,labelsep=sufetwo,skip=6pt,justification=centering}
\captionsetup[figure]{position=bottom,font=sufeten,labelfont=bf,labelsep=sufetwo,skip=6pt,justification=centering}
\renewcommand{\listfigurename}{图\quad 目\quad 录}
\renewcommand{\listtablename}{表\quad 目\quad 录}
\setlength{\heavyrulewidth}{1.0pt}
\setlength{\lightrulewidth}{0.5pt}
\renewcommand{\arraystretch}{1.15}

\ctexset{
  chapter/format={\centering\heiti\zihao{3}},
  chapter/beforeskip={24pt},
  chapter/afterskip={18pt},
  chapter/aftername={\qquad},
  chapter/afterindent=true,
  section/format={\centering\heiti\zihao{4}},
  section/name={第,节},
  section/number={\chinese{section}},
  section/afterindent=true,
  section/beforeskip={24pt},
  section/afterskip={6pt},
  section/aftername={\qquad},
  subsection/format={\heiti\fontsize{13pt}{15.6pt}\selectfont},
  subsection/name={,、},
  subsection/number={\chinese{subsection}},
  subsection/beforeskip={12pt},
  subsection/afterskip={6pt},
  subsection/aftername={},
  chapter/name={第,章},
  chapter/number={\chinese{chapter}},
  contentsname={目\quad 录},
  figurename={图},
  tablename={表}
}
\renewcommand{\thetable}{\thechapter.\arabic{table}}
\renewcommand{\thefigure}{\thechapter.\arabic{figure}}

\let\cleardoublepage\clearpage
\setlength{\parindent}{2em}
\setlength{\parskip}{0pt}
\makeatletter
\@afterindenttrue
\let\@afterindentfalse\@afterindenttrue
\makeatother
\linespread{1.37931}
\raggedbottom
\pagestyle{fancy}
\fancyhf{}
\fancyhead[C]{\zihao{5}\songti \nouppercase{\leftmark}}
\fancyfoot[C]{\zihao{5}\thepage}
\renewcommand{\headrulewidth}{0.4pt}
\renewcommand{\chaptermark}[1]{\markboth{#1}{#1}}
\renewcommand{\sectionmark}[1]{}

\newtheorem{proposition}{命题}[chapter]
\newtheorem{theorem}{定理}[chapter]
\newcommand{\E}{\mathbb{E}}
\newcommand{\indic}{\mathbf{1}}
\newcommand{\method}{FinPlanRAG}
\newcommand{\anchorevolve}{AnchorEvolve}
\newcommand{\sufehead}[1]{\chapter*{#1}\markboth{#1}{#1}\addcontentsline{toc}{chapter}{#1}}
\newcommand{\sufesec}[1]{\par\vspace{12pt}\noindent{\heiti\zihao{4}#1}\par\vspace{6pt}\phantomsection\addcontentsline{toc}{section}{#1}}

% ==================== 高亮修订标记 ====================
% 黄色底纹 = 新增段落/表格行；蓝色文字 = 修改或新增句子。
\newcommand{\rev}[1]{\textcolor{blue}{#1}}
\newcommand{\revhl}[1]{\colorbox{yellow!30}{#1}}
\newtcolorbox{revbox}{
  colback=yellow!12,
  colframe=yellow!55!black,
  breakable,
  boxrule=0pt,
  left=3pt, right=3pt, top=3pt, bottom=3pt,
  before skip=4pt, after skip=4pt
}
\newtcolorbox{oldnote}{
  colback=gray!8,
  colframe=gray!55!black,
  breakable,
  boxrule=0pt,
  left=3pt, right=3pt, top=3pt, bottom=3pt
}

\title{\method：面向金融混合查询的契约路由与锚点演化框架}
% 预答辩/查重匿名版：不写入作者；终稿恢复 \author{高深深}
\author{}
\date{2026年11月}

\begin{document}
\setcounter{tocdepth}{3}
\setcounter{secnumdepth}{3}

% ==================== 封面（预答辩/查重匿名版） ====================
% 按学院通知：封面仅保留论文题目和专业，删除姓名、学号、学院、指导教师等个人信息；
% 原创性声明、致谢、个人简历在本阶段不排入，终稿阶段恢复。
\begin{titlepage}
\thispagestyle{empty}
\noindent
\begin{minipage}[t]{0.5\linewidth}
\vspace{0pt}
\includegraphics[width=2.6cm]{figures/sufe_logo.png}
\end{minipage}%
\begin{minipage}[t]{0.5\linewidth}
\raggedleft
{\zihao{5}\songti 学校代码：10272\par}
\end{minipage}

\vspace{1.1cm}
\begin{center}
{\calligfont\fontsize{40pt}{48pt}\selectfont 上海财经大学\par}
\vspace{0.45cm}
{\zihao{3}\bfseries SHANGHAI UNIVERSITY OF FINANCE AND\par}
{\zihao{3}\bfseries ECONOMICS\par}
\vspace{0.55cm}
{\zihao{2}\heiti 硕士学位论文\par}
\vspace{0.25cm}
{\zihao{3}\bfseries MASTER\quad DISSERTATION\par}
\end{center}

\vspace{2.6cm}
\begin{center}
{\zihao{4}\songti
\renewcommand{\arraystretch}{1.75}
\begin{tabular}{r@{\quad}l}
论文题目 & \underline{\makebox[9.6cm][c]{\textbf{FinPlanRAG：面向金融混合查询的}}} \tabularnewline
 & \underline{\makebox[9.6cm][c]{\textbf{契约路由与锚点演化框架}}} \tabularnewline
专\qquad 业 & \underline{\makebox[9.6cm][c]{金融工程与量化投资}} \tabularnewline
\end{tabular}\par}
\end{center}
\end{titlepage}

% ==================== 原创性声明与版权使用授权书（终稿恢复） ====================
% 按学院通知：预答辩与查重阶段不排入原创性声明和版权使用授权书，终稿阶段恢复。
\frontmatter
\pagestyle{fancy}
% \sufehead{学位论文原创性声明}
% {\zihao{-4}\songti
% \setlength{\parindent}{2em}
% \hspace*{2em}本人郑重声明：所呈交的学位论文，是本人在导师的指导下，独立进行研究工作所取得的成果。除文中已经注明引用的内容外，本论文不含任何其他个人或集体已经发表或撰写过的作品成果。对本文的研究做出重要贡献的个人和集体，均已在文中以明确方式标明。本人完全意识到本声明的法律结果由本人承担。\par
% \vspace{0.7cm}
% \noindent 论文作者签名：\includegraphics[height=1.0cm]{figures/signature.png} \hfill 日期：2026年11月10日\par
% }

% \vspace{0.8cm}
% \phantomsection
% \addcontentsline{toc}{chapter}{学位论文版权使用授权书}
% \begin{center}
% {\zihao{4}\heiti 学位论文版权使用授权书\par}
% \vspace{0.15cm}
% {\zihao{4}\heiti （硕士学位论文用）\par}
% \end{center}
% \vspace{0.15cm}
% {\zihao{-4}\songti
% \setlength{\parindent}{2em}
% \hspace*{2em}本人完全了解上海财经大学关于收集、保存、使用学位论文的规定，即：按照有关要求提交学位论文的印刷本和电子版本。上海财经大学有权保留并向国家有关部门或机构送交本论文的复印件和扫描件，允许论文被查阅和借阅。学校可以将学位论文的全部或部分内容编入有关数据库进行检索和传播，可以采用影印、缩印或扫描等复制手段保存、汇编学位论文。\par
% \vspace{0.7cm}
% \noindent 论文作者签名：\includegraphics[height=1.0cm]{figures/signature.png} \hfill 导师签名：\underline{\hspace{3cm}}\par
% \vspace{0.4cm}
% \noindent 日期：2026年11月10日 \hfill 日期：2026年11月10日\par
% }

% ==================== 修订说明（高亮版） ====================
\clearpage
\sufehead{修订说明}
\begin{revbox}
\noindent\textbf{高亮约定：}黄色底纹表示本次修订新增或重写的段落/表格行；蓝色文字表示修改或新增的句子；未高亮文字为原文保留。
\end{revbox}

\begin{revbox}
\noindent\textbf{叙事重构。}本次修订把论文主线从“router + 两个并列方法”收束为一条契约链：\textbf{契约级形式化}（输出契约作为查询级决策变量）→ \textbf{确定契约控制分支}（义务检索，标准但强，不作为主要算法创新）→ \textbf{前瞻契约主分支}（AnchorEvolve 的可演化锚点证据状态）→ \textbf{全动作损失感知读出}（核心经验贡献）。Router 被明确为契约选择器的简单接口，不再作为第一创新点。
\end{revbox}

\begin{revbox}
\noindent\textbf{算法修订。}(1) 用截止时点 8-K 原文构造带段落编号的文本证据，锚点必须引用 \texttt{P\#}/\texttt{E\#} 来源与日期，公司名默认脱敏；(2) 五类资格判定（vacuous/stale/adversarial/hard/soft）落到确定性代码；(3) 结果揭示后才更新的锚点类型可靠性 Beta 后验与时点先验，构成时间演化；(4) 全动作损失感知读出：在 3 个点标签、7 个非空裁决集和弃权上取最小期望损失，取代原先“点答案 vs 最高密度 $1-\alpha$ 可信集”的受限 gate；(5) 所有 baseline 共用同一个 gate，修正原先的不公平比较。
\end{revbox}

\begin{revbox}
\noindent\textbf{主要新结果（DeepSeek Flash，227 条 XBRL 核验断言）。}旧系统 Fusion 受限 gate 的决策损失为 3.604；全动作读出将其降到 2.863（配对 95\% CI $[-1.031,-0.423]$），GLM 与 Qwen 分别从 3.648/3.577 降到 2.956/2.837。DeepSeek 演化版最终系统（新锚点 + 时间演化 + 演化先验 RF + 全动作读出）为 Acc .551 / Bal .489 / Macro .503 / Loss 2.890，相对旧系统 Loss 差 $-0.714$，CI $[-1.057,-0.352]$；演化先验 RF 单独为 .555/.503/.508/2.828。公平 gate 下 Fusion 相对最强 baseline 点估计更低，Qwen 显著，DeepSeek/GLM 的最强 baseline 差异不显著；准确率相对 RF 的 +2.7 点为点估计（McNemar $p\approx0.69$），不作为显著性主张。Router 结构不变；在 527 条已验证子集上，混合效用由 .389 升到 .437。
\end{revbox}

\begin{revbox}
\noindent\textbf{需要读者注意。}附录 A 的十一项指标和原表 6.2/6.3 是\textbf{旧受限 gate} 口径；第 6 章与附录 G 给出全动作口径与演化版结果。原稿的类型化义务章节保留为设计描述；主实验的确定分支实现是确定性词项覆盖，完整类型化 planner 在 \texttt{src/finplan\_rag/planner.py} 中作为扩展。
\end{revbox}

% ==================== 中文摘要 ====================
\clearpage
\phantomsection
\addcontentsline{toc}{chapter}{摘\texorpdfstring{\quad}{ }要}
\begin{center}
{\zihao{-2}\heiti \method：面向金融混合查询的\\[0.2cm]
契约路由与锚点演化框架\par}
\vspace{0.2cm}
{\zihao{-2}\heiti 摘\quad 要\par}
\end{center}
\vspace{0.2cm}
\markboth{摘\quad 要}{摘\quad 要}
{\zihao{-4}\songti\sloppy
\setlength{\parindent}{0pt}
\setlength{\emergencystretch}{2em}
\hspace*{2em}\rev{金融文本分析同时包含确定信息问题与前瞻性断言。前者要求从截止时点以前的申报文件中读取或计算已陈述事实，后者要求判断尚未实现的未来结果。现有检索增强生成、断言验证与选择性预测通常预先固定输出接口，因而容易把前瞻性断言压缩为无保留的点答案，或把本可闭合的事实问题强制为分类裁决。本文把输出契约表述为查询级决策变量：系统先判断事实是否已在截止时点存在，再选择闭合答案契约或点答案/裁决集/弃权契约。Router 是这一契约选择器的简单接口，不作为本文的主要算法创新。确定契约由义务检索控制分支实现；前瞻契约由 AnchorEvolve 实现：它从截止时点 8-K 原文抽取带来源的 L2 锚点，经五类资格判定、时点可靠性演化与信念更新得到可审计证据状态，再由全动作损失感知读出转换为点答案、裁决集或弃权。}
\hspace*{2em}\rev{本文在 DocFinQA、227 条 XBRL 核验收入指引断言和 FinDVer 上，使用 DeepSeek Flash、GLM 5.3 Flash 和 Qwen3-8B 与六个具名基线进行统一评估。全动作损失感知读出在三个基座上显著降低决策损失：Fusion 从 3.604 降到 2.863、从 3.648 降到 2.956、从 3.577 降到 2.837（配对区间均排除零）。DeepSeek 演化版最终系统达到准确率 .551、平衡准确率 .489、宏 F1 .503、决策损失 2.890，相对旧系统显著降损（−0.714，CI [−1.057,−0.352]）；演化先验 RF 诊断达到 .555/.503/.508/2.828。公平 gate 下 Fusion 的点估计低于全部具名基线，Qwen 上显著；DeepSeek/GLM 上相对最强基线的差异不显著。结果表明：契约形式化负责组织问题，锚点演化负责维护可审计证据状态，全动作校准读出负责控制最终风险。}
\noindent{\zihao{-4}\heiti 关键词：}{\zihao{-4}\songti 金融问答\hspace{0.6em}检索增强生成\hspace{0.6em}前瞻性断言\hspace{0.6em}契约路由\hspace{0.6em}锚点演化\hspace{0.6em}全动作决策\hspace{0.6em}校准}
}

% ==================== 英文摘要 ====================
\clearpage
\phantomsection
\addcontentsline{toc}{chapter}{Abstract}
\begin{center}{\arial\bfseries\zihao{-2} Abstract\par}\end{center}
\vspace{0.6cm}
\markboth{Abstract}{Abstract}
{\arial\zihao{-4}\selectfont\sloppy
\setlength{\parindent}{0pt}
\setlength{\emergencystretch}{3em}
\hspace*{2em}\rev{Financial text analysis combines determinate questions, whose answers are stated or computable at a cutoff, with forward-looking claims, whose outcomes have not yet been realized. Existing retrieval-augmented generation, claim-verification, and selective-prediction systems usually fix the output interface in advance. We formulate the output contract as a query-level decision: the system first decides whether the fact exists at the cutoff, then selects either a closed-answer contract or a point/credible-set/abstention contract. The router is a deliberately simple interface to this contract selector, not the main algorithmic contribution. The determinate contract is implemented by an obligation-retrieval control branch. The forward-looking contract is implemented by AnchorEvolve: it extracts source-cited L2 anchors from the point-in-time 8-K excerpt, applies a five-type qualification, evolves anchor-type reliability with revealed outcomes, updates a belief, and converts the resulting auditable evidence state into a point answer, a loss-optimal adjudication set, or an abstention through a full-action loss-aware readout.}
\hspace*{2em}\rev{We evaluate FinPlanRAG on DocFinQA, 227 XBRL-verified revenue-guidance claims, and FinDVer with DeepSeek Flash, GLM 5.3 Flash, and Qwen3-8B against six named baselines. The full-action readout significantly lowers decision loss on all three backbones: Fusion decreases from 3.604 to 2.863, from 3.648 to 2.956, and from 3.577 to 2.837, with paired intervals excluding zero. On DeepSeek, the evolved system reaches .551 accuracy, .489 balanced accuracy, .503 macro-F1, and 2.890 decision loss, a significant reduction from the old system (-0.714, CI [-1.057,-0.352]); the evolved-prior RF diagnostic reaches .555/.503/.508/2.828. Under the same gate, Fusion is lower in point estimate than every named baseline and is significant on Qwen; the differences against the strongest baselines on DeepSeek and GLM are not significant. The results support a division of labor in which the contract formulation organizes the problem, anchor evolution maintains an auditable evidence state, and the full-action calibrated readout controls final risk.}
}
\vspace{0.4cm}
{\sloppy\noindent{\arial\bfseries\zihao{-4} Key Words:} {\arial\zihao{-4}\selectfont financial QA\hspace{0.8em}RAG\hspace{0.8em}forward-looking claims\hspace{0.8em}contract routing\hspace{0.8em}calibrated credible sets}\par}

% ==================== 目录 ====================
\clearpage
\tableofcontents
\markboth{目\quad 录}{目\quad 录}
\clearpage
\phantomsection
\addcontentsline{toc}{chapter}{图\texorpdfstring{\quad}{ }目\texorpdfstring{\quad}{ }录}
\listoffigures
\markboth{图\quad 目\quad 录}{图\quad 目\quad 录}
\clearpage
\phantomsection
\addcontentsline{toc}{chapter}{表\texorpdfstring{\quad}{ }目\texorpdfstring{\quad}{ }录}
\listoftables
\markboth{表\quad 目\quad 录}{表\quad 目\quad 录}
\clearpage

\mainmatter
\zihao{-4}
\setlength{\parindent}{2em}
\zihao{-4}

% ==================== 第1章 ====================
\chapter{绪论}
\label{ch:intro}

\section{研究背景与问题提出}

金融文本分析连接财务报告理解、事实检索与业绩判断，但真实查询并不服从单一的问答契约。确定信息问题要求系统从截止时点以前的10-K、10-Q、8-K及其表格、脚注和XBRL事实中读取或计算已经陈述的结果，评价重点是答案准确性与证据可定位性；前瞻性断言则涉及尚未实现的营收指引、产能投放、监管审批或合同履约，其真值在截止时点之后才会揭示，合理输出可能是点判断、可信集或弃权。管理层预测研究也表明，指引同时包含区间、时点和可验证性，不能等同于已实现业绩\citep{hirst2008management}。围绕长文档、表格推理和证据获取，金融问答与检索增强生成已经形成系统方法\citep{lewis2020rag,chen2021finqa,zhu2021tatqa,islam2023financebench,reddy2024docfinqa}，主动检索、自适应检索和智能体框架进一步优化了检索时机与证据使用\citep{jiang2023flare,jeong2024adaptiverag,yao2023react,shinn2023reflexion}；金融断言验证研究将陈述映射为支持、反驳或证据不足\citep{zhao2024findver,rangapur2025finfact}，选择性预测则通过拒识控制低置信度输出\citep{chow1970optimum,geifman2017selective}。这些研究分别改进了证据获取、断言核验和风险控制，却通常预先固定回答接口。固定契约用于混合负载时，确定问题可能被压缩为分类裁决，前瞻性断言也可能被阅读器输出为无保留的点答案。尚未被充分处理的关键问题因此不是单一检索器或分类器的性能，而是系统能否先判断一条查询应遵循事实检索契约还是前瞻性判断契约，并在相应契约下维护证据状态。本文据此提出FinPlanRAG：给定查询$q$与截止时点$c$，LLM路由器选择义务检索或AnchorEvolve；前者以类型化义务驱动时点取证并返回带证据编号的闭合答案，后者以稳定关系、行业模式和实例三层锚点组织前瞻性证据，依据可靠性、作用域和时效更新结果信念。点答案、可信集或弃权由校准读出生成，使契约选择、证据演化与风险控制在同一时点约束下形成可审计链条。

\section{研究内容与方法}

本文按照查询契约识别、分支证据建模和统一实证检验展开。首先，将金融混合查询形式化为确定信息问题与前瞻性断言问题的并集，并以截止时点界定可用信息，所有检索、特征构造、校准和读出选择均遵守时点约束，以控制前视偏差\citep{lopezdeprado2018advances}。其次，构建由LLM路由器、义务检索和AnchorEvolve组成的总体框架。路由器仅依据查询与截止时点元数据返回契约标签、置信度和理由；确定信息问题被解析为类型化义务，系统围绕未覆盖义务迭代检索，由LLM阅读器判断证据闭包并生成带编号答案；前瞻性断言被映射到有限结果空间，AnchorEvolve抽取L0稳定关系、L1行业模式和L2实例锚点，通过五类资格判定更新信念，校准读出据此生成点答案、可信集或弃权。最后，本文在DocFinQA、XBRL核验收入指引和FinDVer上，使用DeepSeek Flash、GLM 5.3 Flash和Qwen3-8B进行同基座比较，并分别检验路由效用、答案质量、证据成本、锚点增量信息、概率校准和非对称决策损失。这样的研究设计使契约识别、证据组织与风险控制能够被分层验证，同时保留端到端系统在混合负载上的整体评价。

\section{主要贡献}

\begin{revbox}
\noindent\textbf{贡献一：契约级形式化。}本文把输出契约从预先固定的接口变成查询级决策变量：系统先判断“事实是否已经在截止时点存在”，再选择闭合答案契约或点答案/裁决集/弃权契约。现有检索增强生成、断言验证和选择性预测分别固定了输出接口，本文把契约选择前置到检索与生成之前，使两类查询共享同一时点约束而保持不同输出形式。Router 是这一契约选择器的\textbf{简单接口}，本文不主张 router 网络结构或分类器设计上的创新；确定分支是第二契约的\textbf{控制分支}，也不是本文的新检索算法。
\end{revbox}

\begin{revbox}
\noindent\textbf{贡献二：可演化的锚点证据状态。}本文提出 AnchorEvolve 的金融断言实例化：从截止时点的 8-K 原文抽取带来源编号的 L2 锚点，以 L0 关系词表和 L1 行业锚点类型组织语义，经过五类资格判定（vacuous/stale/adversarial/hard/soft）形成可审计状态；结果公开揭示后，锚点类型可靠性以 Beta 后验更新，时点先验与演化可靠性进入信念更新。该状态的价值是\textbf{可审计、时点正确、可作为决策层输入}，而不是单独承担点预测；本文不主张 LLM 锚点后验单独优于强基线。
\end{revbox}

\begin{revbox}
\noindent\textbf{贡献三：全动作损失感知读出与公平评价。}原先的读出只在点答案与最高密度 $1-\alpha$ 可信集之间选择；本文枚举 3 个点标签、全部 7 个非空裁决集和弃权，取期望损失最小的动作。该读出在同一后验和同一损失下是严格改进，并在 DeepSeek、GLM、Qwen 三个基座上显著降低决策损失。本文同时把所有具名基线放到同一个读出契约下重算，修正了“只有 FinPlanRAG 能使用集合/弃权”的不公平比较；公平 gate 下 Fusion 的点估计低于全部具名基线，Qwen 显著，DeepSeek/GLM 上相对最强基线的差异不显著。
\end{revbox}

\begin{revbox}
\noindent\textbf{本文不主张的内容。}(1) 不主张 router 结构创新：它是一个零样本 LLM 二分类接口；(2) 不主张确定分支是新检索算法：它是标准但强的控制分支，主实验使用确定性词项义务覆盖；(3) 不主张 LLM 锚点后验单独胜过 RF：演化先验 RF 仍是强诊断，锚点层的价值在审计与语义状态；(4) 不主张准确率相对 RF 的差异显著：+2.7 点是点估计，损失相对旧系统的下降才是显著结果。
\end{revbox}

% ==================== 第2章 ====================
\chapter{文献综述}
\label{ch:related}

本章围绕金融问答与断言验证、检索增强生成与自适应检索、选择性预测与可信集、锚点约束与时点证据四条线索梳理相关研究，并在此基础上归纳研究缺口与关键挑战。

\noindent\textbf{一、金融问答与断言验证}

金融问答基准测试模型在申报文件、表格和报告上的数值推理能力。FinQA研究基于财务报表的数值推理\citep{chen2021finqa}，TAT-QA研究表格与文本的混合推理\citep{zhu2021tatqa}，FinanceBench研究真实金融文档上的开放问答\citep{islam2023financebench}，FinBen研究金融领域的综合评测\citep{xie2024finben}。DocFinQA在FinQA问题的基础上加入完整10-K上下文，把证据选择从段落级扩展到长文档级\citep{reddy2024docfinqa}。这些基准的共同点是答案在文档中已经存在，评价指标主要是精确匹配或词级F1。它们衡量的是阅读器能否在给定文档中找到并计算答案，而不是系统能否判断查询是否应当以闭合答案回应。

金融申报文件的结构与通用问答语料不同。10-K和10-Q通常包含多级标题、跨页表格、脚注和附注，同一指标可能在正文、表格和管理层讨论中以不同口径出现；8-K则按事项编号披露业绩、指引、并购和监管事件。XBRL数据提供结构化财务事实，但季度口径、单位、重述和可用日期仍需与文本对齐。通用阅读器擅长语义匹配，却不一定能处理期间、单位和证据闭合，因此金融问答需要把文档结构、时点约束和数值计算纳入同一证据策略。

断言验证基准测试模型判断一条陈述是否被证据支持、反驳或无法判断。FinDVer面向金融领域的断言验证，要求模型在长文档和混合内容中给出裁决和解释\citep{zhao2024findver}；Fin-Fact面向证据驱动的金融事实核查，强调多模态证据和解释生成\citep{rangapur2025finfact}。这类任务的输出是离散裁决，评价指标是裁决准确率和解释质量。验证器通常假设陈述已经被给定，且证据已经被检索，因而没有回答“这条陈述是否应当进入验证流程”的问题。

近期金融RAG系统进一步加入类型化证据、期间约束、来源追踪和语料对齐检索\citep{mansouri2026finrank,jiang2026finrate,chen2026hcrag,chi2026finsagent,zhou2026fincards}。这些系统在证据表示和检索配方上明显优于通用RAG，但输出接口仍通常是固定答案或固定裁决。它们改善了到达阅读器的证据，却没有改变阅读器与验证器之间的契约边界。对于同时包含事实问题和前瞻性断言的金融分析场景，这种固定接口仍然是主要限制。

\noindent\textbf{二、检索增强生成与自适应检索}

RAG将检索器与生成器组合，假设答案可以由检索到的上下文拼装\citep{lewis2020rag}。主动检索和自适应检索进一步决定何时检索或检索多少上下文\citep{jiang2023flare,jeong2024adaptiverag}，智能体框架则交错进行推理、检索和自我修正\citep{yao2023react,shinn2023reflexion}。Self-RAG通过自我反思决定是否检索和如何评价证据\citep{asai2024selfrag}，CRAG在检索质量不足时触发纠正性检索\citep{yan2024crag}，FLARE在生成过程中按置信度触发检索\citep{jiang2023flare}，Adaptive-RAG按问题复杂度选择检索策略\citep{jeong2024adaptiverag}。

这些方法改善了检索过程和证据使用，但输出通常仍是生成的答案。它们可以决定“是否检索”和“检索多少”，却不能决定“返回答案、可信集还是弃权”。在确定信息问题上，这种输出是合适的；在前瞻性断言上，单一生成答案无法表达结果分布、证据作用域和决策风险。检索策略的自适应并不等于输出契约的自适应，二者是两个不同层次的决策。

本文的义务检索分支借用了自适应检索的迭代思想，但把停止条件从“检索轮数”改为“义务覆盖与阅读器充分性”。这一改动使证据预算服务于问题的语义结构，而不是服务于固定的$k$。AnchorEvolve则把输出契约从生成答案扩展为点答案、可信集和弃权的组合，使系统可以在证据不足时保留不确定性。

\noindent\textbf{三、选择性预测、校准与可信集}

拒识选项在识别错误和拒识率之间进行权衡\citep{chow1970optimum}，选择性分类通过置信度阈值在低置信度输入上弃权\citep{geifman2017selective}。现代分类器通常存在校准误差，温度缩放是常用基线\citep{guo2017calibration}；代价敏感学习为不同错误赋予不同价格\citep{elkan2001cost}；分裂共形预测则把得分转换为在可交换性下具有分布无关覆盖率的集合\citep{vovk2005algorithmic,angelopoulos2021gentle}。分类共形方法进一步构造自适应集合\citep{romano2020classification}，分布漂移下的自适应共形方法则允许覆盖率随时间调整\citep{gibbs2021adaptive}。

这些方法提供了损失、置信度和集合构造的工具，但通常假设任务和回答接口已经固定。选择性分类器在“答案”和“弃权”之间选择，验证器在“支持”“反驳”和“证据不足”之间选择，共形方法在给定标签空间上构造集合。它们没有解决“标签空间本身是否适用于当前查询”的问题。对于前瞻性断言，标签空间是结果桶，证据是锚点，输出还需要说明证据的作用域和时效；这些结构不在标准选择性预测的范围内。

本文借用选择性预测的损失框架和共形集合的覆盖思想，但把标签空间、证据资格和输出动作放在同一个锚点状态上。校准随机森林和温度缩放用于控制后验概率，损失感知门用于在答案、可信集和弃权之间选择。与标准共形方法相比，本文的可信集由后验质量构造，覆盖率在实验中被报告，但并未假设交换性；这是金融时点数据与标准共形设定之间的差异。

\noindent\textbf{四、锚点、部分识别与时点证据}

部分识别表示与现有证据一致的结果集合\citep{manski2003partial}。在金融机器学习中，模型决策只能使用当时可得的信息，否则会产生前视偏差\citep{lopezdeprado2018advances}。本文借用约束和可行集的思想，将锚点定义为来源与假设之间的关系：硬锚点可以排除某个结果桶，软锚点只重新加权结果信念；证据过滤和模式过滤均按截止时点执行。

与部分识别相比，本文的锚点不是完整的结构模型，而是由LLM从文本中抽取的语义约束。锚点带有类型、取值、作用域、来源、时间、可靠性和有效期，因而可以被逐条检查。部分识别给出的是在给定模型下与证据一致的参数集合，锚点给出的是在给定断言下与证据一致的结果桶约束。两者都强调“证据不足时不要强行点识别”，但本文更关注证据的可审计性和可更新性。

时点约束贯穿本文的两个分支。义务检索只接受可用日期不晚于截止时点的文档，AnchorEvolve只接受截止时点之前的证据，随机森林只使用截止时点之前可得的特征，校准和权重选择只在开发集与验证集上进行。这一约束使“当时可得信息”成为系统设计的一部分，而不是事后评价的附加条件。

\noindent\textbf{五、研究缺口与关键挑战}

现有研究分别推进了长文档检索、金融问答、断言验证和低置信度拒识，但尚未把查询契约选择与前瞻性证据演化纳入同一框架。本文认为，这一缺口可以分解为三个相互衔接的挑战。三个挑战分别对应查询层、证据层和决策层，也构成本文框架的设计依据。

第一个挑战是查询契约的识别。金融查询的表面形式与真实契约并不一一对应。“预期股利”“未来支付金额”等表述包含未来时态，但答案可能在文件中已经给出；管理层指引包含精确数字，却是关于尚未实现结果的前瞻性断言。系统如果按关键词或时态路由，会把确定问题送入裁决分支，或把前瞻性断言送入阅读器。真正的判断依据是查询所要求的事实是否在截止时点已经存在，以及输出应当采用闭合答案还是可校准判断。这一挑战的难点不在分类器本身，而在于契约是查询、证据与时点之间的关系，而不是查询文本的主题标签。

第二个挑战是前瞻性证据的状态化。现有验证器通常把断言直接映射为支持、反驳或证据不足，阅读器则直接生成答案。两者都没有保留证据的作用域、可靠性、时效和来源，也没有在证据变化时更新判断。对于前瞻性断言，证据往往只支持某个结果桶、只在某个时间范围内有效、只对某个公司或行业成立；把这些信息压缩为一个标签，系统就无法解释判断为何改变，也无法在结果揭示后更新可靠性。本文认为，前瞻性断言需要一个可审计的中间状态，使证据到信念、信念到决策的每一步都可以追溯。

第三个挑战是语义判断与风险校准的分工。LLM能够从文本中抽取锚点、解释证据和形成语义判断，但其概率输出往往存在类别偏置和校准误差；统计模型可以使用结构化特征进行概率校准，却缺少对文本证据的理解。若让LLM直接输出最终裁决，系统可能过度自信；若只用统计模型，系统又无法利用文本中的语义约束。本文认为，合理的分工是让LLM构造语义状态，让校准程序控制决策风险，并在验证集上选择读出方式，而不是在测试集上比较后择优。

上述三个挑战对应本文的三项贡献。契约识别由查询契约路由回应；证据状态化由AnchorEvolve的三层锚点、五类资格判定和信念更新，以及义务检索回应；语义与校准的分工由校准后端、损失感知门和可信集读出回应。三项贡献分别对应查询层、证据层和决策层，构成逐层回应。

% ==================== 第3章 ====================
\chapter{研究设计：研究问题、假设与总体框架}
\label{ch:framework}

\section{研究问题}

金融文本分析中的查询可以分为两类。确定信息问题要求返回申报文件在截止时点已经陈述或可以计算的事实，答案通常是数值、段落或明确的预测。前瞻性断言问题要求判断尚未实现的未来结果，其真值取决于截止时点之后披露的信息，正确输出可能是一个点答案、一个可信集，或者在证据不足时弃权。设查询为$q$，截止时点为$c$，工作负载是确定问题集合$\mathcal Q_D$与前瞻性断言集合$\mathcal Q_V$的混合。系统调用LLM路由器
\begin{equation}
r(q,c)\in\{D,V\}
\end{equation}
判断查询属于哪一类，再将确定查询送入义务检索分支，将前瞻性断言送入锚点演化分支。

两个分支返回不同的输出形式。义务检索分支返回闭合答案$a$、置信度与证据编号，证据不足时返回\texttt{insufficient}。AnchorEvolve返回点答案$h^*$、可信集$R_\alpha$、动作$\pi$与信念$b$，动作可以是回答、可信集裁决或弃权。路由器只读取原始查询和截止时点元数据，不读取“这是断言”的任务标签，也不读取测试结果。

上述问题可以展开为三个研究问题：
\begin{enumerate}[label=RQ\arabic*.,leftmargin=2em,itemsep=2pt,topsep=2pt]
\item 契约识别：仅依据查询与截止时点元数据，系统能否在混合负载上区分确定信息契约与前瞻性断言契约？
\item 确定信息取证：类型化义务与覆盖驱动的证据获取，能否在更低证据成本下提高确定信息问答的准确率与覆盖率？
\item 前瞻性裁决：三层锚点状态与校准读出，能否在保持分类准确率的同时降低决策损失，并提供可信集或弃权？
\end{enumerate}
三个问题分别对应查询层、证据层和决策层。RQ1检验路由器能否识别契约；RQ2检验义务检索能否以更低成本提高确定问答的准确率和覆盖率；RQ3检验锚点状态与校准读出能否把语义信息转换为低风险决策。

\section{研究假设}

三个研究问题对应三个可检验假设。假设的表述同时给出检验口径和否证条件，使后文实验能够直接检验框架是否有效，而不是只描述框架结构。

\begin{revbox}
\noindent\textbf{H1（契约可识别性）}：仅依据查询与截止时点元数据，简单 LLM 接口能够区分确定信息契约与前瞻性断言契约；在较强基座上，路由混合效用接近 oracle，并高于 Always-RAG 与 Always-Anchor。检验口径为 573 条构造混合负载上的路由准确率、混淆矩阵和混合效用。若路由器在多数基座上低于较好的固定分支，则 H1 不成立。\textbf{边界：}H1 检验的是“契约选择是否可学”，不是 router 结构创新；本文不主张 router 分类器设计优于已有自适应检索路由。
\end{revbox}

\begin{revbox}
\noindent\textbf{H2（确定契约的控制结果，非主创新）}：确定的词项义务覆盖与阅读器闭包，能够在低于固定 16 段的平均证据成本下提高 DocFinQA 精确匹配和覆盖率，并迁移到 FinDVer。检验口径为 DocFinQA 的 300 题固定样本和 FinDVer testmini。若分支在多数基座上不高于固定 16 段阅读器，或证据成本没有下降，则 H2 不成立。\textbf{边界：}H2 是控制分支的可信度检查，用来保证 router 的另一半不是 strawman；本文不主张新的检索算法，也不把类型化元数据过滤作为主实验已实现的性质。
\end{revbox}

\begin{revbox}
\noindent\textbf{H3a（全动作损失感知读出）}：在同一个后验和同一个非对称损失下，枚举 3 个点标签、7 个非空裁决集和弃权的 full-action gate，决策损失低于“点答案 vs 最高密度 $1-\alpha$ 可信集”的受限 gate。检验口径为 227 条 XBRL 核验断言上的决策损失，三个基座分别配对检验。若 full-action 读出在多数基座上不低于受限 gate，则 H3a 不成立。
\end{revbox}

\begin{revbox}
\noindent\textbf{H3b（演化锚点状态与融合）}：文本锚点 + 时间演化 + 演化先验 RF 的最终系统，在保持可审计状态的同时，决策损失不劣于静态先验/静态可靠性和 RF 诊断。检验口径为 DeepSeek 上的 227 条测试断言、validation 选择的融合权重、以及 static/evolved 与 new/old 锚点消融。若演化系统在多数比较中劣于静态或 RF 诊断，则 H3b 不成立。\textbf{边界：}H3b 的准确率/平衡准确率/宏 F1 只报告点估计；现有 227 条样本上相对 RF 的准确率差异不显著，不能写成“显著优于 RF”。
\end{revbox}

本文的假设层级是：H1 是接口可行性，H2 是控制分支可信度，H3a 是核心决策层贡献，H3b 是演化状态的增量检验。第\ref{ch:empirical}章按这个层级组织实验；H3a 在当前三个基座上成立，H3b 的损失增量是点估计且不显著，论文如实报告这一边界。


\section{总体框架}

FinPlanRAG由三个模块组成。路由器是一个独立的LLM调用，负责判断查询契约；义务检索分支负责确定信息问题；AnchorEvolve分支负责前瞻性断言。路由器与两个分支之间通过结构化输出连接，路由器不读取分支输出，分支不改变路由结果。这样的结构使路由错误和分支错误可以被分开诊断：路由错误表现为查询被送入不匹配的契约，分支错误表现为在正确契约下的证据或判断错误。

路由器读取查询文本与截止时点元数据，返回路由标签、置信度和简短理由。提示词定义两个契约：确定信息请求要求返回申报文件在截止时点陈述或可计算的事实、数值、段落、估计或预测；前瞻性断言要求判断尚未实现的未来结果，例如实际营收是否落在指引区间内、产品是否按时发布、目标是否达成。提示词明确说明管理层指引是前瞻性断言，不是已实现事实，即使其中包含数字。

路由器不读取已实现结果，也不读取两个分支的输出。它只看到查询和截止时点元数据，因此在部署时可以按查询缓存，成本是一次短LLM调用。路由器的错误类型也具有语义含义：确定问题被送到AnchorEvolve，通常是因为查询文本包含未来时态或预期措辞；前瞻性断言被送到义务检索，通常是因为断言中包含可检索的数字或期间。这些错误为改进路由器提供了可解释的信号。

两个分支共享路由器，但不共享输出契约。义务检索分支返回闭合答案，AnchorEvolve返回点答案、可信集或弃权。共享路由器使系统只需一次契约判断；不共享输出契约使每个分支可以针对自己的证据结构和风险结构设计。图\ref{fig:framework}给出总体结构，第\ref{ch:determinate}章和第\ref{ch:anchor}章分别给出两个分支的完整设计。

\begin{revbox}

\noindent\textbf{接口而非创新。}Router 在这里是一个简单的零样本 LLM 接口，用来把契约选择显式化；本文不主张 router 网络结构或分类器设计上的新意。本文的算法贡献集中在两处：确定契约的控制分支（标准但强的义务检索），以及前瞻契约的 AnchorEvolve——文本锚点、五类资格、时间演化的可靠性/先验，和全动作损失感知读出。
\end{revbox}

\begin{figure}[H]
\centering
\resizebox{0.98\linewidth}{!}{%
\begin{tikzpicture}[
  font=\normalsize,
  >=Latex,
  box/.style={draw, rounded corners=3pt, align=center, text width=2.6cm, minimum height=0.9cm, inner sep=3pt, fill=white},
  panel/.style={draw, rounded corners=6pt, inner sep=7pt, fill=gray!5},
  arrow/.style={->, thick, draw=black!65},
  lbl/.style={font=\footnotesize, text=black!70}
]
\node[box, fill=gray!10] (q) at (0,0) {查询 $q$\\截止时点 $c$};
\node[box, fill=gray!18, text width=2.6cm] (r) at (2.7,0) {LLM 路由器\\契约判断};

\node[box] (o1) at (5.6,2.0) {类型化义务};
\node[box] (o2) at (8.6,2.0) {时点检索};
\node[box] (o3) at (11.6,2.0) {LLM 阅读器};

\node[box] (a1) at (5.6,-2.0) {三层锚点};
\node[box] (a2) at (8.6,-2.0) {资格与信念};
\node[box] (a3) at (11.6,-2.0) {校准与损失门};

\begin{scope}[on background layer]
\node[panel, fit=(o1)(o3), label={[font=\bfseries]above:义务检索分支}] {};
\node[panel, fit=(a1)(a3), label={[font=\bfseries]below:AnchorEvolve 分支}] {};
\end{scope}

\node[box, text width=2.1cm] (ans) at (14.6,2.0) {闭合答案};
\node[box, text width=2.1cm] (pt) at (14.6,0.6) {点答案};
\node[box, text width=2.1cm] (cs) at (14.6,-0.6) {可信集};
\node[box, text width=2.1cm] (ab) at (14.6,-2.0) {弃权};
\node[lbl, above] at (14.6,2.65) {输出};

\draw[arrow] (q) -- (r);
\draw[arrow] (r.east) -- ++(0.7,0) |- (o1.west);
\draw[arrow] (r.east) -- ++(0.7,0) |- (a1.west);
\draw[arrow] (o1) -- (o2);
\draw[arrow] (o2) -- (o3);
\draw[arrow] (a1) -- (a2);
\draw[arrow] (a2) -- (a3);
\draw[arrow] (o3) -- (ans);
\draw[arrow] (a3) -- (pt);
\draw[arrow] (a3) -- (cs);
\draw[arrow] (a3) -- (ab);

\draw[dashed, draw=black!45] (0.9,3.0) -- (0.9,-3.0);
\node[lbl, above] at (0.9,3.05) {时点约束 $V_c$};
\end{tikzpicture}
}
\caption{FinPlanRAG总体框架。路由器把查询送入义务检索分支或AnchorEvolve；前者返回闭合答案，后者返回点答案、可信集或弃权。左侧虚线表示时点约束。}
\label{fig:framework}
\end{figure}


\section{时点约束与评价协议}

系统在截止时点$c$形成判断，可用信息由可采纳证据集合$V_c$给出。设证据$e$的可用时间为$a(e)$，则
\begin{equation}
V_c=\{e:a(e)\le c,\ e\text{在}c\text{时未被撤回}\}.
\end{equation}
策略$\pi$在$c$时有效，当且仅当它是$\sigma(V_c)$可测的。这一条件排除了使用未来结果、未来披露或未来修订的证据。路由器不读取结果，锚点抽取器不接收结果，随机森林只使用截止时点之前可得的特征，校准和权重选择只在开发集与验证集上进行。测试结果不参与任何拟合或选择。

评价口径与输出契约对应。确定信息问题使用精确匹配、词级F1、答案覆盖率和证据成本；前瞻性断言使用准确率、平衡准确率、宏F1、MCC、AUROC、AUPRC、ECE、Brier、决策损失、可信集覆盖率和平均集合大小；混合负载使用路由准确率、混淆矩阵和混合效用。所有任务共享同一时点约束和同一套测试集不参与选择的原则，因此不同分支的结果可以在同一框架下比较。

评价协议由三部分组成。第一，主结果使用配对自助法给出置信区间，准确率差异使用McNemar检验；所有区间在断言或问题级别重采样。第二，随机森林使用五个折种子，测试年份的训练集来自其他公司折，且结果日期早于测试年份1月1日；测试公司不出现在训练集中。第三，所有LLM响应按请求负载的哈希缓存，提示词、运行脚本、逐条预测和表格生成脚本随论文发布，API密钥只从环境变量读取。

这一协议的目的是让每个假设都有对应的证据边界。H1对应573条混合负载；H2对应DocFinQA的300题固定样本和FinDVer testmini；H3a/H3b对应227条XBRL核验断言、全动作读出消融、锚点消融和演化消融。

% ==================== 第4章 ====================
\chapter{义务检索分支：Obligation Retrieval}

\begin{revbox}

\noindent\textbf{角色与边界（修订）。}本章的确定分支是第二契约的\textbf{控制分支}，不是本文的主要算法创新。主实验实现的是确定性的词项义务覆盖（数字/月份/期间词项 + BM25 + 阅读器闭包 + 置信度回退）；4.1 的类型化元组与 metadata 一致性是设计目标，完整实现在 \texttt{src/finplan\_rag/planner.py} 与附录 F。主文只报告控制结果，逐条件、token 成本与 planner-16 对照放在附录与局限。
\end{revbox}
\label{ch:determinate}

义务检索分支可以概括为四步：把问题表示为类型化义务，构建时点证据空间，按未覆盖义务迭代取证，由阅读器生成带证据编号的闭合答案。四步依次对应本章四节。

\begin{revbox}

\noindent\textbf{实现口径。}第 4.1--4.4 节保留类型化义务的设计描述；当前主结果的实现是 \texttt{obligation\_terms} 的词项覆盖版本，chunk 的 \texttt{available\_at} 为固定时间戳，未实现文档版本/元数据过滤。这一差异在局限与附录 F 中明确说明，避免把设计目标写成已实现的实验性质。
\end{revbox}

\section{状态表示：类型化义务}

义务检索分支把自然语言问题绑定为可检索、可核对的类型化义务。义务元组定义为
\begin{equation}
o=(e,f,y,p,m,u),
\end{equation}
其中$e$为实体，$f$为表单或文档类型，$y$为财年，$p$为财期，$m$为指标，$u$为单位。以“某公司2023财年第四季度营业收入是多少”为例，实体是公司，表单是10-K或10-Q，财年是2023，财期是第四季度，指标是营业收入，单位是美元。义务项从问题和查询元数据中抽取；抽取器只负责识别约束，不负责回答。

文档只有在可用日期不晚于截止时点、且元数据与义务一致时才被接受。元数据一致性检查覆盖实体、表单、财年、财期、指标和单位。与只按问题文本检索相比，义务元组把实体、期间和单位变成可检查的约束。例如，问题文本中的“2023年”可能指财年、自然年或季度，BM25无法区分；义务元组把“2023财年第四季度”固定为可核对的期间，从而减少跨期文档和单位不一致带来的错误。

类型化义务还提供了证据覆盖的度量。设问题$q$的义务项集合为$T(q)$，已获取证据集合为$A_t$，则未覆盖义务项为
\begin{equation}
U_t=T(q)\setminus\{t\in T(q):\exists a\in A_t,\ t\in a\}.
\end{equation}
未覆盖项既用于构造下一轮检索查询，也用于判断何时停止。义务覆盖不是答案正确性的充分条件，但它把“证据是否覆盖了问题的关键约束”变成可观察的状态，使系统可以在证据不足时继续检索或返回证据不足，而不是在固定轮数后强行生成答案。

\section{证据空间：时点语料与块索引}

分支把文档切分为段落块和表格块。段落块按行合并，单块不超过1200字符；表格块按连续表格行合并，单块不超过4000字符。每个块记录文档编号、可用时间、文本、词项和元数据。切块的目标是保留表格行与表头的对应关系，同时避免单块过长导致阅读器输入超出预算。对于跨页表格，分支保留表头并在块内重复列名，以减少阅读器对列含义的误判。

分支使用BM25建立索引，按问题与未覆盖义务项检索候选块。BM25的优势是无需训练、可复现、对长文档稳定；缺点是不理解语义，因此需要义务覆盖和阅读器反馈补足。检索查询由原问题与$U_t$拼接而成，候选块按得分排序，去重后最多新增两个块。这样做的原因是：一次加入过多块会稀释阅读器的注意力，也会增加token成本；一次加入过少块则可能需要更多轮检索。两块的设置是在证据覆盖和调用成本之间的折中。

时点过滤在检索前执行。只有可用日期不晚于截止时点的文档进入索引；被后续版本撤回或替代的文档被排除。对于同一文档的多个版本，分支保留截止时点之前的最新版本。时点过滤使系统不会使用未来披露的信息，也避免了“文档存在但当时尚未发布”的前视偏差。

\section{取证策略：覆盖驱动的迭代检索}

第$t$轮检索后，分支把新增块加入证据集$A_t$，调用LLM阅读器判断义务是否闭合。停止条件有三类：阅读器报告证据充分且$U_t=\varnothing$；阅读器给出非空答案、置信度不低于0.8且$U_t=\varnothing$；预算$B$耗尽。默认$B=8$，并报告$B=12$和$B=16$。前两类条件把停止权交给证据状态和阅读器判断，第三类条件提供成本上限。

覆盖驱动的取证把停止变量从检索前的固定预算移到义务覆盖和阅读器充分性上。对于简单问题，分支可以在两三轮内闭合；对于复杂问题，分支可以继续检索直到预算耗尽。这种自适应性使平均证据成本低于固定预算，同时在难题上保留足够的证据预算。

置信度门控是覆盖驱动取证的补充。当阅读器置信度低于0.8或报告证据不足时，分支回退到固定16段证据集重新阅读。回退不是对每个问题支付16段成本，而是对规划器无法闭合的问题支付更高成本。实验表明，回退只在规划器无法闭合的问题上触发，DeepSeek Flash上约为24\%，GLM 5.3 Flash上比例更高；它弥补了规划器单独的准确率缺口，同时保持平均段落数低于固定16段。完整流程见算法\ref{alg:obligation}。



\section{答案生成：阅读器与闭包}

阅读器接收问题与已获取块，返回结构化答案
\begin{equation}
(\text{answer},\text{confidence},\text{sufficient},\text{missing},\text{evidence\_ids},\text{rationale}).
\end{equation}
答案要求为短字符串；若证据不足，返回\texttt{insufficient}并将\texttt{sufficient}置为\texttt{false}。分支把阅读器的\texttt{missing}字段拼入下一轮查询，形成阅读器驱动的取证。例如，阅读器报告缺少“第四季度”或“单位”，下一轮查询会加入这些词项，从而把检索引向缺失的证据。

闭包规则要求答案附带证据编号。每个答案都可以追溯到具体段落或表格块，证据不足时有明确的缺失信息。这一规则使义务检索分支的输出可以被下游审计：读者可以检查答案是否由证据支持，也可以检查系统在证据不足时是否选择弃权而不是猜测。

解析失败在评价中计为错误。阅读器被要求返回JSON，但基座可能返回解释性文字或格式错误的JSON。分支对解析失败不进行静默修复，而是记录为错误，以避免评价口径被解析器放宽。Qwen3-8B在DocFinQA上返回较多解释性文字，导致精确匹配接近地板；这一结果作为基座能力边界报告，不用于义务检索分支的主比较。

\begin{algorithm}[H]
\caption{义务检索的单轮取证}
\label{alg:obligation}
\begin{algorithmic}[1]
\Require 问题$q$，截止时点$c$，义务项$T(q)$，预算$B$
\Ensure 答案$a$或\texttt{insufficient}
\State $A\gets\varnothing$; $U\gets T(q)$
\For{$t=1,\dots,B$}
  \State $C_t\gets\textsc{BM25}(q\cup U,\,V_c)$
  \State $A\gets A\cup\textsc{Top2}(C_t\setminus A)$
  \State $(a,\text{conf},\text{suff},\text{miss},\text{ids})\gets\textsc{Reader}(q,A)$
  \State $U\gets T(q)\setminus\textsc{Covered}(A)$
  \If{$\text{suff}=\text{true}$ and $U=\varnothing$}
    \State \Return $a$
  \EndIf
  \If{$\text{conf}\ge0.8$ and $a\ne\varnothing$ and $U=\varnothing$}
    \State \Return $a$
  \EndIf
  \State $q\gets q\cup\text{miss}$
\EndFor
\State $A\gets\textsc{Fixed16}(q,V_c)$
\State $(a,\text{conf},\text{suff},\dots)\gets\textsc{Reader}(q,A)$
\If{$\text{suff}=\text{true}$}
  \State \Return $a$
\Else
  \State \Return \texttt{insufficient}
\EndIf
\end{algorithmic}
\end{algorithm}

义务检索的时间成本主要由BM25检索和阅读器调用组成。BM25在本地索引上完成，不产生API成本；阅读器调用次数由停止条件决定，最坏情况下为预算$B$加上一次固定16段回退。设文档块数为$N$，单次BM25检索的复杂度为$O(N\log N)$，相对于LLM调用可以忽略。

可审计性来自证据编号和缺失信息。每个答案都附带实际使用的段落或表格块编号，证据不足时返回具体的缺失项。审计者可以检查答案是否由证据支持，也可以检查系统在证据不足时是否选择弃权而不是猜测。

% ==================== 第5章（重构版） ====================
% ==================== 第5章（重构版） ====================
\chapter{AnchorEvolve：前瞻性断言分支}
\label{ch:anchor}

\begin{revbox}
\noindent\textbf{本章修订说明。}本章把 AnchorEvolve 从“一次性锚点抽取 + 静态后验”重构为“时点文本证据 → 五类资格 → 时间演化可靠性/先验 → 信念 → 校准融合 → 全动作损失感知读出”的链条。高亮部分为新增或重写内容；未高亮的概念定义保留原稿。实现边界在 5.3、5.4 和 5.6 节明确标注，避免把设计目标写成已实现的实验性质。
\end{revbox}

AnchorEvolve 可以概括为六步：定义假设空间，构造文本锚点，过滤锚点，演化可靠性并更新信念，校准融合，最后在全动作空间上读出。前两步对应状态表示与状态构造，第三步对应五类资格判定，第四步对应时间演化与信念更新，第五步对应校准融合，第六步对应全动作损失感知读出。

\section{状态表示：假设空间与三层锚点}

前瞻性断言的真值取决于截止时点之后披露的信息。系统在截止时点只能形成关于未来结果的判断，不能把判断当作已实现事实。AnchorEvolve 把断言表示为有限假设空间上的信念。通用设计中，假设空间是时机与阶段的笛卡尔积：
\begin{equation}
\mathcal H=\{(t,s):t\in\mathcal T,\ s\in\mathcal S\},
\end{equation}
其中$\mathcal T$为季度或年度时机，$\mathcal S$为试验、量产等阶段。对于收入指引，本文实例化为
\begin{equation}
\mathcal H=\{\text{低于区间},\text{区间内},\text{高于区间}\}.
\end{equation}
断言自带的时机与阶段作为弱先验，不作为事实。系统不接收已实现结果，只接收截止时点之前的证据。

AnchorEvolve 的核心状态是三层锚点。L0 是稳定关系本体：界、先后、前置时间、持续时间、守恒、比率、单位换算和作用域。L0 不随行业变化，为锚点提供可复用的关系类型。例如，“当前季度值落在指引区间内”是一个界关系，“新产品在指引季度之后投产”是一个先后关系，“产能爬坡需要两个季度”是一个前置时间关系。L0 的作用是限制锚点可以表达的关系种类，使不同行业的锚点共享同一套语义骨架。

L1 是行业模式：每个锚点类型是（名称、L0 关系、字段、生效轮次）四元组。收入指引的初始模式包含指引宽度、当前季度增长、公司历史、行业历史、文本信号和其他。L1 把 L0 关系绑定到具体字段和生效规则。例如，“指引宽度”绑定到界关系，字段是上限与下限之差除以中点；“当前季度增长”绑定到比率关系，字段是当前季度值相对指引中点的比例。L1 带有生效轮次，使新发现的锚点类型不会在同一轮中使用。

L2 是锚点实例：编号、类型、取值或区间、单位、作用域、来源、时间、可靠性和有效期。L2 是 LLM 抽取器的直接输出，也是信念更新的输入。每个 L2 锚点都记录来源和时间，因此可以被逐条检查。

\begin{revbox}
\noindent\textbf{表示边界。}L0/L1/L2 在当前实现中是“关系词表 + 类型 schema + 实例”的三层表示，不是三个独立训练的模块：L1 类型用于可靠性演化，L0 关系用于语义分类，L2 实例进入信念更新。L1 的文本信号类型（\texttt{text\_signal}）是本次修订新增的类型，用于承载 8-K 原文中的管理层评论、分部/需求信号；其余五类沿用原稿。
\end{revbox}

\begin{figure}[htbp]
\centering
\begin{tikzpicture}[
  font=\small,
  >=Latex,
  box/.style={draw, rounded corners=2pt, align=center, text width=3.4cm, minimum height=1.05cm, inner sep=2pt},
  small/.style={draw, rounded corners=2pt, align=center, text width=3.0cm, minimum height=1.0cm, inner sep=2pt},
  arrow/.style={->, thick, draw=black!70}
]
\node[box] (q) at (0,0) {前瞻性断言 $q$\\截止时点 $c$};
\node[box] (e) at (0,-1.8) {时点 8-K 原文\\段落 $P\#$ / 证据 $E\#$};
\node[box, fill=gray!8] (x) at (0,-3.6) {LLM 锚点抽取\\L0/L1/L2};
\node[small] (h) at (-3.8,-5.5) {硬锚点\\可靠性 $\ge$ 阈值};
\node[small] (s) at (0,-5.5) {软锚点\\似然更新};
\node[small] (f) at (3.8,-5.5) {空泛/过期/对抗\\过滤};
\node[box] (b) at (0,-7.2) {结果桶信念 $b$};
\node[box] (ev) at (-3.8,-8.9) {时间演化状态\\类型可靠性 / 时点先验};
\node[box] (ba) at (0,-8.9) {锚点后验 $b_A$};
\node[box, fill=gray!8] (br) at (3.8,-8.9) {结构化 RF $b_R$};
\node[box] (sel) at (0,-10.6) {校准融合 $b_F$};
\node[box] (out) at (0,-12.3) {全动作损失感知读出\\点答案 / 裁决集 / 弃权};
\draw[arrow] (q) -- (e);
\draw[arrow] (e) -- (x);
\draw[arrow] (x) -- (h);
\draw[arrow] (x) -- (s);
\draw[arrow] (x) -- (f);
\draw[arrow] (h) -- (b);
\draw[arrow] (s) -- (b);
\draw[arrow] (b) -- (ba);
\draw[arrow] (ev) -- (ba);
\draw[arrow] (ba) -- (sel);
\draw[arrow] (br) -- (sel);
\draw[arrow] (sel) -- (out);
\end{tikzpicture}
\caption{AnchorEvolve 修订版机制。时点 8-K 原文经 LLM 抽取 L0/L1/L2 锚点；五类资格判定将硬/软锚点送入信念更新，空泛/过期/对抗锚点被过滤；时间演化状态（类型可靠性 Beta 后验与时点先验）参与信念更新；锚点后验与结构化 RF 融合后，由全动作损失感知读出生成点答案、裁决集或弃权。}
\label{fig:anchormechanism}
\end{figure}

\section{状态构造：时点文本证据与锚点抽取}

\begin{revbox}
\noindent\textbf{文本证据。}每条 asserted claim 的输入是 \texttt{claim\_source.url} 指向的、截止时点可得的 8-K exhibit，而不是结果揭示的 earnings release。系统保留开头若干段落（标题、当季实际、下一季指引），再按 revenue/guidance/outlook/segment/demand 等关键词选取段落；每个段落带稳定编号 $P1,P2,\dots$，结构化证据带 $E1,E2,\dots$。公司名默认脱敏，prompt 只接收匿名公司标识、行业、截止时点、目标期间、指引区间、当前季度值和时点先验。
\end{revbox}

锚点抽取器是一个 LLM 调用。输入包括前瞻性断言、公司匿名标识、行业、截止时点、目标期间、指引区间、当前季度值与可靠性、观测先验和行业先验、演化类型可靠性、可用证据项和 8-K 段落摘录。输出包括一组 L2 锚点以及三类桶的分数。每个锚点记录类型、关系、取值、作用域、来源编号、日期、可靠性、硬软标记、隐含桶和理由；三类分数给出低于区间、区间内和高于区间的边际判断。

\begin{revbox}
\noindent\textbf{提示词约束。}修订版 prompt 明确要求：（1）锚点必须引用 \texttt{P\#} 或 \texttt{E\#} 来源，并给出不晚于 cutoff 的日期；（2）\texttt{reliability} 是“隐含桶正确的校准概率”，不是推理置信度；（3）指引区间定义三个桶，但区间本身不是“结果落在区间内”的证据；（4）证据不足时返回 \texttt{missing}，不编造锚点。提示词与结构化输出 schema 见附录 D 的修订版。
\end{revbox}

抽取器的三类分数是边际信号。分支把它与锚点更新后的信念等权混合，得到锚点后验。混合后的后验保留锚点的结构化信息，同时利用抽取器对三类结果的边际判断；混合权重 $\gamma$ 是配置项，默认 $0.5$。

\section{状态过滤：五类资格判定}

每个锚点在更新信念之前经过资格判定。修订版把五类判定写成确定性规则；所有阈值在配置文件中给出，不在代码中写死。

\begin{table}[H]
\centering
\zihao{5}
\caption{五类锚点资格判定。}
\label{tab:qualification}
\begin{tabular}{p{2.2cm}p{5.2cm}p{4.8cm}}
\toprule
类型 & 判定条件 & 处理 \\
\midrule
\rowcolor{yellow!18} vacuous & 无隐含桶，或锚点不提供有效桶信息 & 不进入更新 \\
\rowcolor{yellow!18} stale & 锚点日期晚于 cutoff，或来源日期晚于 cutoff & 不进入更新 \\
\rowcolor{yellow!18} adversarial & 引用的段落/证据不在时点摘录中，或作用域不匹配 & 不进入更新 \\
\rowcolor{yellow!18} hard & 显式约束且可靠性不低于硬锚点阈值 & 限制可行桶（当前实现为点质量+小底） \\
\rowcolor{yellow!18} soft & 其余有效锚点 & 以似然更新信念 \\
\bottomrule
\end{tabular}
\end{table}

\begin{revbox}
\noindent\textbf{实现边界。}当前代码对 hard 锚点采用“隐含桶点质量 + 小底”的更新；多个 hard 锚点的可行集交集、冲突检测和回退到软更新是论文原设计，代码中尚未实现，作为扩展接口保留。空泛/过期/对抗三类已经实现确定性检查；在当前收入指引数据上主要触发 hard/soft，这与原稿的边界说明一致。
\end{revbox}

五类判定分开处理证据的不同维度。硬、软刻画约束强度；空泛刻画信息量；对抗刻画作用域；过期刻画时效。把它们合并为一个可靠性分数会丢失诊断信息：一条高可靠性但作用域错误的锚点不应进入更新，一条作用域正确但信息量低的锚点也不应主导信念。分开判定使每个被拒绝的锚点都有明确理由，也使审计者可以检查过滤规则是否合理。

\section{状态更新：时间演化与信念}

\begin{revbox}
\noindent\textbf{时间演化。}对 cutoff 为 $c$ 的断言，系统只使用结果公开揭示日早于 $c$ 的锚点和结果。揭示日取 \texttt{earnings\_events.json} 的 earnings event 日期；缺失时回退到 10-Q 日期；若揭示日不早于 $c$，该结果在 $c$ 时不可见。这样，可靠性更新和先验更新都是时点正确的，而不是把未来结果泄漏到当前判断。
\end{revbox}

信念从时点先验 $b_0$ 开始。修订版不再使用固定的开发集先验，而是用历史已揭示结果的平滑计数：
\begin{equation}
b_0(k)=\frac{n_k+1}{n+3},\qquad k\in\mathcal H,
\end{equation}
其中 $n_k$ 是在 $c$ 之前已公开揭示的、结果为桶 $k$ 的断言数，$n=\sum_k n_k$。当 $n=0$ 时退化为均匀先验。

每个锚点类型 $t$ 维护一个 Beta 可靠性后验。设截至 $c$ 已揭示的、类型 $t$ 的锚点中，隐含桶与结果一致的数量为 $m_t$，总数为 $n_t$，则
\begin{equation}
r_t=\frac{m_t+\alpha}{n_t+\alpha+\beta},
\end{equation}
其中 $\alpha,\beta$ 是配置中的 Beta 先验参数。修订版把演化可靠性与原 LLM 主观可靠性做收缩：
\begin{equation}
\rho_t=\lambda r_t+(1-\lambda)r_{\mathrm{LLM}},
\end{equation}
其中 $\lambda$ 是收缩系数（默认 $0.5$）。$\rho_t$ 作为软锚点更新的指数。

软锚点的似然为
\begin{equation}
L_e(k)=
\begin{cases}
p_e, & k=j,\\
(1-p_e)/(|\mathcal H|-1), & k\ne j,
\end{cases}
\end{equation}
其中 $j$ 为隐含桶，$p_e$ 是配置项（默认 $0.7$）。软锚点执行
\begin{equation}
b(k)\propto b(k)\,L_e(k)^{\rho_t}.
\end{equation}
硬锚点把信念置为隐含桶上的点质量加一个小底，以避免后续软锚点无法恢复。每一步更新后对信念做归一化。

抽取器的三类分数是边际信号。分支把它与锚点更新后的信念按权重 $\gamma$ 混合：
\begin{equation}
b_A=\gamma\,b+(1-\gamma)\,s,
\end{equation}
其中 $s$ 为三类分数的归一化向量，$\gamma$ 为配置项（默认 $0.5$）。

锚点后验的优点是语义可解释：每个桶的概率变化可以追溯到具体锚点，硬锚点的排除和软锚点的加权都有记录。缺点是概率质量不一定校准；LLM 抽取器可能系统性地高估某一类桶，锚点可靠性是抽取器给出的主观量，不是经过验证的频率。时间演化用已揭示结果校准类型可靠性，但并不能消除所有类别偏置；因此第 5.5 节仍引入结构化 RF 和校准融合。

\section{决策读出：校准融合与全动作损失感知门}

\begin{revbox}
\noindent\textbf{从受限 gate 到全动作 gate。}原稿的读出只在“点答案”和“最高密度 $1-\alpha$ 可信集”之间比较。修订版把动作空间扩展为：3 个点标签 $a\in\mathcal H$、全部 7 个非空裁决集 $R\subseteq\mathcal H$、以及弃权。对每个动作计算期望损失并取最小。该读出在同一个后验和同一个损失下是严格改进；原受限 gate 是新动作空间的一个子集。
\end{revbox}

校准后端是一个随机森林，输入为截止时点可得的结构化特征：指引宽度、当前季度值相对指引中点的比例、指引与当前季度的比值、历史结果计数与频率、前次指引增长与宽度、目标季度、截止年份、行业，以及本次修订加入的演化先验特征（peer/sector 历史、公司 surprise 历史、日历、当前季度相对 low/high 的位置、可靠性标志）。随机森林在开发集与验证集上训练，并在公司不重叠折上做温度缩放；测试公司不出现在训练集中。

锚点后验与校准随机森林后验 $b_R$ 融合：
\begin{equation}
b_F=w\,b_A+(1-w)\,b_R,
\end{equation}
其中 $w$ 只在验证集上从配置的权重网格中选择，测试集不参与选择。修订版不再把 $w$ 写死为 $0.1$；DeepSeek 上新锚点 + 演化先验 RF 的验证选择为 $w=0.20$。

给定输出后验 $b^\ast$，系统在点回答、裁决集和弃权之间选择期望损失最低的动作。回答损失为
\begin{equation}
L_{\text{ans}}(a,y)=
\begin{cases}
0, & a=y,\\
\ell_{\text{over}}, & a=\text{高于},y=\text{低于},\\
\ell_{\text{under}}, & a=\text{低于},y=\text{高于},\\
\ell_{\text{wrong}}, & \text{其他}.
\end{cases}
\end{equation}
集合损失为
\begin{equation}
L_{\text{set}}(R,y)=\mathbf{1}[y\notin R]\,\ell_{\text{miss}}
+\ell_{\text{width}}\max(0,|R|-1),
\end{equation}
弃权损失为$\ell_{\text{abstain}}$。全动作读出在 $b^\ast$ 下计算三类动作的期望损失并取最小：
\begin{equation}
\pi^\star=\arg\min_{\pi\in\{\text{answer},\text{adjudicate},\text{abstain}\}}
\mathbb E_{y\sim b^\ast}\big[L(\pi,y)\big].
\end{equation}
默认参数为$\ell_{\text{over}}=12$、$\ell_{\text{under}}=8$、$\ell_{\text{wrong}}=10$、$\ell_{\text{miss}}=8$、$\ell_{\text{width}}=2$、$\ell_{\text{abstain}}=5$、$\alpha=0.1$。

\begin{revbox}
\noindent\textbf{可信集与裁决集的区分。}$1-\alpha$ 可信集
\begin{equation}
R_\alpha=\arg\min_{R\subseteq\mathcal H}\{|R|:\textstyle\sum_{k\in R}b^\ast(k)\ge 1-\alpha\}
\end{equation}
仍作为覆盖率和审计口径报告，但它不再等同于决策集。全动作 gate 选择的裁决集 $R^\star$ 是损失最优集，可以在 $1-\alpha$ 集更宽时选择更窄的集合，或在不确定性高时选择弃权。DeepSeek 上 $1-\alpha$ 集的平均大小约 2.8、覆盖率约 .99；损失最优动作集的平均大小约 1.9、覆盖率约 .86。这是给定非对称损失下的最优权衡，不应把“覆盖率低于名义值”误读为缺陷。
\end{revbox}

损失参数的选择影响动作分布。本文使用一组固定参数报告主结果，并在附录中给出参数含义与敏感性分析的计划；参数不是从测试集调出来的，而是对应金融分析中对过度乐观、漏报和集合宽度的相对定价。不同机构可以根据自身风险偏好替换参数，而不改变框架结构。

\section{时点演化与算法}

AnchorEvolve 在结果揭示后更新来源可靠性和锚点类型可靠性。对每个使用过的证据项，分支记录其与已实现桶是否一致，更新来源可靠性、锚点类型可靠性和分类器权重。第 $r$ 轮的结果在第 $r$ 轮全部断言裁决完成后才更新状态，并从第 $r+1$ 轮开始使用。新发现的锚点类型生效轮次为 $r+1$，因此触发发现的断言不能使用该锚点。该规则保证自演化是时点的。收入指引实验报告第一轮；多轮循环作为扩展接口保留。

\begin{algorithm}[H]
\caption{AnchorEvolve 单条断言的演化与全动作裁决}
\label{alg:anchor}
\begin{algorithmic}[1]
\Require 断言$q$，截止时点$c$，时点 8-K 分段$E_c$，演化状态$S_c$，结构化 RF 后验$b_R$
\Ensure 动作$\pi$，点答案$h^\ast$，裁决集$R^\ast$，后验$b_F$
\State $A\gets\textsc{ExtractAnchors}(q,c,E_c,S_c)$
\State $A_q\gets\textsc{Qualify}(A,q,c)$ \Comment{vacuous/stale/adversarial/hard/soft}
\State $b_A\gets\textsc{BeliefUpdate}(S_c,A_q)$
\State $b_F\gets w\,b_A+(1-w)\,b_R$
\State $(\pi,h^\ast,R^\ast)\gets\textsc{FullActionGate}(b_F)$
\State $S_{c^+}\gets\textsc{RevealUpdate}(S_c,c)$ \Comment{仅用揭示日早于下一 cutoff 的结果}
\State \Return $(\pi,h^\ast,R^\ast,b_F)$
\end{algorithmic}
\end{algorithm}

算法的时间成本主要由锚点抽取的输入和输出长度决定。单条断言的输入与输出通常在数千 token 量级；随机森林推断在本地完成，不产生 API 成本。与端到端微调相比，该流程不需要锚点级人工标注，锚点类型和资格规则可以逐条检查；与直接提示模型输出类别相比，AnchorEvolve 保留从证据到锚点、从锚点到信念、从信念到动作的中间状态。

\begin{revbox}
\noindent\textbf{可审计性与边界。}每个被采纳或拒绝的锚点都有来源编号、时间、作用域和理由；每个桶的信念变化都可以追溯到具体锚点；每个动作都有对应的期望损失。需要强调：当前实现尚未做人工审计实验，因此“可审计”是设计属性而非已测量的用户收益；LLM 抽取器仍可能出错，错误的硬锚点可能排除正确桶，可靠性阈值和冲突检查只能缓解、不能消除。多轮演化、来源可靠性单独更新和硬锚点可行集交集是后续工作。
\end{revbox}


% ==================== 第6-7章与附录（重构版） ====================
% ==================== 第6章（重构版） ====================
\chapter{实证结果与分析}
\label{ch:empirical}

\begin{revbox}
\noindent\textbf{本章修订说明。}本章按修订后的主线重排结果：先报告 router 的接口可行性（6.3），再把确定分支作为第二契约的控制结果报告（6.4），然后把\textbf{全动作损失感知读出}作为核心经验结果，在三个基座上与受限 gate 和公平 gated 基线比较（6.5）；6.6--6.8 报告 DeepSeek 上的文本锚点、时间演化和演化先验结果。原稿表 6.2/6.3 与附录 A 的十一项指标是\textbf{旧受限 gate} 口径；本章表格使用全动作口径，旧表仅作历史对照。
\end{revbox}

\section{数据与任务}

本文在三个数据集和一个混合负载上评估 FinPlanRAG。三个数据集分别对应长文档确定问答、前瞻性断言裁决和外部迁移，混合负载用于检验路由。这样的设计使框架的不同主张分别接受检验，而不是用一个数据集的平均结果替代全部证据。

DocFinQA 在 FinQA 问题的基础上加入完整 10-K 上下文\citep{reddy2024docfinqa}。本文使用测试集中的 300 题固定样本，答案使用发布的金标准。FinQA 程序常量只作为证据召回率的银标准，不作为阅读器输入。该数据集的挑战在于证据分散在长文档中，表格行与段落需要联合阅读，阅读器必须在大量无关内容中定位分子、分母和期间。

收入指引基准包含 2010 年至 2025 年美国 SEC 8-K 文件中的 1,005 条季度营收指引断言。每条断言包含指引区间、截止时点、当前季度值和证据项。实际季度营收来自 SEC XBRL，XBRL 没有独立季度时使用披露文件回退。标签为低于区间、区间内、高于区间。数据按时间划分为 529 条开发集（2010 年至 2018 年）、203 条验证集（2019 年至 2021 年）和 273 条测试集（2022 年至 2025 年）。测试集包含 126 条高于、110 条区间内和 37 条低于。主结果使用 227 条具有 XBRL 核验或 XBRL 覆盖标签的测试断言；46 条纯文本标签单独用于敏感性检验。该数据集的挑战在于标签来自自动核验而非人工标注，类别不平衡明显，且前瞻性断言的真值在截止时点尚未实现。

FinDVer testmini 包含 700 条带财务报告的断言验证样本\citep{zhao2024findver}，作为外部迁移集。本文不在 FinDVer 上调参，只把在 DocFinQA 上设计的义务检索策略迁移过去，以检验证据策略是否依赖特定数据分布。

混合负载由 300 条 DocFinQA 问题和 273 条收入指引断言组成，共 573 条查询。确定查询的效用为词级 F1；前瞻性查询的效用为 $1-L/10$，其中 $L$ 为决策损失。错误路由记 0 效用。Oracle 路由器把每条查询送到正确分支；Always-RAG 把所有查询送到义务检索分支；Always-Anchor 把所有查询送到 AnchorEvolve。混合效用把两类任务的分数放在同一尺度上，使路由策略可以与固定分支策略比较。

\section{基座模型、具名基线与评价指标}

本文使用三套基座：DeepSeek Flash、GLM 5.3 Flash 和 Qwen3-8B\citep{qwen2025qwen3}。Qwen3-8B 在本地部署。所有方法使用同一基座和相同的提示模板，以比较证据策略和输出契约的差异。

义务检索分支与固定 BM25 预算 1、2、4、8、16 段、Self-RAG-style、CRAG-style、FLARE-style、Adaptive-RAG-style 以及义务检索规划器比较。四个具名 RAG 方法根据原论文描述和公开接口进行同基座适配，保留各自的检索或反思策略，使用与义务检索分支相同的段落块、阅读器和评分，因此差异来自证据策略。

AnchorEvolve 与六个具名方法比较：FinDVer-style\citep{zhao2024findver}、Fin-Fact-style\citep{rangapur2025finfact}、Self-RAG-style\citep{asai2024selfrag}、CRAG-style\citep{yan2024crag}、FLARE-style\citep{jiang2023flare} 和 Adaptive-RAG-style\citep{jeong2024adaptiverag}。每个方法根据原论文描述和公开接口进行同基座适配，通过同一个三分类适配层运行。适配层把原始输出映射为低于、区间内、高于，并提供相同的时点证据块和相同的六条历史断言。学习组件不可用或不兼容时，用使用官方判据的 API 调用替代，并记录替代关系。六个条件保留各自原论文的判据，不是原系统在其原始数据上的未修改端到端运行。

\begin{revbox}
\noindent\textbf{公平评价协议（修订）。}原稿只给 FinPlanRAG 使用集合/弃权动作，而六个具名基线按点标签计算损失；这会把“是否有 gate”误读为“方法是否更好”。修订版把同一个全动作损失感知读出应用到所有方法：每个方法的 label/confidence/credible\_set 都转换为后验，并在 3 个点标签、7 个非空裁决集和弃权上取最小期望损失。点准确率仍按各自 argmax 报告；损失、覆盖率、集合大小按统一动作报告。
\end{revbox}

评价指标分为三组。DocFinQA 报告精确匹配、词级 F1、答案覆盖率、平均获取段落数、平均输入 token 和银标准证据召回率。收入指引报告准确率、平衡准确率、宏 F1、MCC、AUROC、AUPRC、ECE、Brier、决策损失、可信集覆盖率和平均集合大小。平衡准确率是各类召回率的平均\citep{brodersen2010balanced,sokolova2009performance}。决策损失使用第\ref{ch:anchor}章参数。损失、ECE 和 Brier 越低越好。置信区间由断言级配对自助法给出，准确率差异使用 McNemar 检验。路由报告路由准确率、混淆矩阵和混合效用。

评价协议在三个层面控制信息泄漏。第一，收入指引按时间划分，测试年份为 2022 年至 2025 年，验证集为 2019 年至 2021 年，开发集为 2010 年至 2018 年。第二，随机森林使用五个折种子；对每个测试年份 $Y$，训练集来自其他公司折，且结果日期早于 $Y$ 年 1 月 1 日；测试公司不出现在训练集中。温度缩放在公司不重叠的内部折上拟合。第三，融合权重 $w$ 只在验证集上选择，候选为配置中的权重网格；测试集不参与拟合、校准、权重选择或组合选择。

所有 LLM 响应按请求负载的哈希缓存。API 密钥只从环境变量读取，不写入结果文件。提示词、运行脚本、逐条预测和表格生成脚本随论文发布。随机森林的折种子、校准方式和融合权重选择在正文中给出。公开结果包含逐条预测、路由决策和汇总指标；LLM 响应缓存不随论文提交，但相同输入可以用相同提示词复现。

\section{路由与混合负载}

\begin{table}[htbp]
\centering
\zihao{5}
\setlength{\tabcolsep}{3pt}
\caption{混合负载路由结果。效用定义见第\ref{ch:empirical}章。}
\label{tab:router}
\resizebox{\linewidth}{!}{%
\begin{tabular}{lrrrrr}
\toprule
基座 & 路由准确率 & 路由器 & 始终RAG & 始终Anchor & Oracle \\
\midrule
DeepSeek & .995 & .403 & .121 & .282 & .403 \\
GLM 5.3 & .990 & .419 & .130 & .292 & .422 \\
Qwen3-8B & .867 & .263 & .018 & .298 & .316 \\
\bottomrule
\end{tabular}
}
\end{table}


表\ref{tab:router}报告 573 条混合负载。DeepSeek Flash 路由器正确路由 570 条（0.995）：300 条确定问题中 297 条，273 条前瞻性断言全部正确。GLM 5.3 Flash 路由器正确路由 567 条（0.990）：297 条确定问题、270 条前瞻性断言。Qwen3-8B 正确路由 271 条确定问题和 226 条前瞻性断言（0.867），其中 29 条确定问题被送到 AnchorEvolve，47 条前瞻性断言被送到义务检索，显示 8B 基座对契约边界的识别较不稳定。

DeepSeek Flash 的混合效用为 0.403，与 oracle 相同，高于 Always-RAG 的 0.121 和 Always-Anchor 的 0.282。GLM 5.3 Flash 的混合效用为 0.419，接近 oracle 的 0.422，并恢复 99.3\% 的 oracle 效用。Qwen3-8B 的混合效用为 0.263，低于 Always-Anchor 的 0.298。结果表明，较大基座能够把确定问答与前瞻性裁决送入相应接口，小基座则需要更强的路由监督或外部信号。

路由错误主要出现在带有未来时态的确定问题，例如“未来支付金额”和“预期股利”。这些问题的答案已经在截止时点文件中披露，但表面表达接近前瞻性断言。该现象说明契约识别依赖事实是否已经存在，而不能仅依据时态或关键词。固定分支策略只能对全部查询采用同一接口，路由器则把契约选择显式化，使路由错误与分支错误可以分别诊断。

\begin{revbox}
\noindent\textbf{接口定位与端到端检查。}Router 是契约选择器的简单接口；本节检验的是“契约选择是否可学”，不是 router 结构创新。修订版进一步用新的前瞻分支损失替换旧损失，在 527 条已验证子集（300 条确定 + 227 条前瞻）上保持 router 预测不变：混合效用从 .389 升到 .437，Always-Anchor 从 .258 升到 .306，router 来自正确确定路由的 .131 优势保持不变。该检查排除了 46 条纯文本标签，因此与 573 条表不完全可比。
\end{revbox}

\section{义务检索控制分支}

\begin{revbox}
\noindent\textbf{控制分支的定位。}本节把确定分支作为第二契约的控制结果报告，而不是本文的主要算法创新。主实验使用确定性的词项义务覆盖（数字/月份/期间词项 + BM25 + 阅读器闭包 + 置信度回退）；第 4 章的类型化元组与 metadata 过滤是设计目标，完整实现在附录 F 的扩展 planner 中。本节只报告控制分支是否足以让 router 的另一半可信；逐条件、token 成本和 planner-16 对照放在附录与局限。
\end{revbox}

\begin{table}[htbp]
\centering
\zihao{5}
\setlength{\tabcolsep}{2pt}
\caption{DocFinQA长文档问答结果。Acc.为精确匹配，Macro F1为宏F1，Paras为平均段落，Cov.为覆盖率。}
\label{tab:docfinqa}
\begin{tabular}{lrrrrrrrr}
\toprule
 & \multicolumn{4}{c}{DeepSeek Flash} & \multicolumn{4}{c}{GLM 5.3 Flash} \\
\cmidrule(lr){2-5}\cmidrule(lr){6-9}
Condition & Acc. & Macro F1 & Paras & Cov. & Acc. & Macro F1 & Paras & Cov. \\
\midrule
Fixed-4 & .333 & .192 & 4.00 & .557 & .400 & .199 & 4.00 & .653 \\
Fixed-8 & .367 & .223 & 8.00 & .603 & .440 & .225 & 8.00 & .720 \\
Fixed-16 & .390 & .234 & 16.00 & .640 & .480 & .249 & 16.00 & .757 \\
Planner B=8 & .417 & .230 & 4.56 & .707 & .450 & .248 & 4.15 & .763 \\
Planner B=16 & .460 & .276 & 6.85 & .773 & .470 & .243 & 5.88 & .840 \\
Self-RAG-style & .373 & .223 & 8.03 & .613 & .447 & .231 & 8.16 & .717 \\
CRAG-style & .373 & .222 & 8.05 & .613 & .447 & .235 & 8.32 & .727 \\
FLARE-style & .330 & .197 & 4.00 & .557 & .393 & .193 & 4.19 & .643 \\
Adaptive-RAG-style & .363 & .211 & 7.22 & .570 & .393 & .198 & 6.51 & .657 \\
\midrule
\textbf{FinPlanRAG (obligation + fallback)} & \textbf{.457} & \textbf{.276} & 10.69 & \textbf{.770} & \textbf{.497} & \textbf{.252} & 10.31 & \textbf{.810} \\
\bottomrule
\end{tabular}
\end{table}


表\ref{tab:docfinqa}报告 DocFinQA 的 300 题固定样本。最终义务检索分支使用 16 段预算，并在阅读器不充分或置信度低于 0.8 时回退到固定 16 段证据集。在 DeepSeek Flash 上，该分支取得 0.457 精确匹配，平均使用 10.69 段；固定 16 段阅读器取得 0.390，平均 16 段；固定 8 段阅读器取得 0.367，平均 8 段。无回退的预算 8 规划器取得 0.417，平均 4.56 段；预算 12 取得 0.427，平均 5.82 段；预算 16 取得 0.460，平均 6.85 段。

在 GLM 5.3 Flash 上，该分支取得 0.497，平均 10.31 段；固定 16 段阅读器取得 0.480，平均 16 段；固定 8 段阅读器取得 0.440，平均 8 段。无回退的预算 8 规划器取得 0.450，平均 4.15 段；预算 12 取得 0.473，平均 5.11 段；预算 16 取得 0.470，平均 5.88 段。回退分支比固定 16 段阅读器高 0.017，平均少用 5.69 段。两个基座上的答案覆盖率分别为 0.770 和 0.810，高于固定预算的 0.640 和 0.757，说明分支更少弃权。

\begin{revbox}
\noindent\textbf{成本与支配关系。}最终分支相对固定 16 段节省段落数，但 DeepSeek 上的平均总 token 为 1712.83，固定 16 段为 1752.50，实际只省 2.3\%；GLM 上为 2632.11 对 3715.89，省 29\%。同时，planner B=16 在 DeepSeek 上是 .460 / 6.85 段 / 1275.31 token，在准确率和 token 上都支配最终分支；最终分支在 GLM 上更强。因此“回退弥补规划器缺口”只对 B=8 成立，对 B=16 不成立；本节把最终分支写成成本-准确率折中，而不是一致占优。
\end{revbox}

\begin{table}[htbp]
\centering
\zihao{5}
\setlength{\tabcolsep}{3pt}
\caption{FinDVer testmini迁移结果（默认阅读器）。}
\label{tab:findver}
\begin{tabular}{llrrr}
\toprule
基座 & 条件 & 准确率 & 宏F1 & 平均段落 \\
\midrule
\multicolumn{5}{l}{\textit{DeepSeek}} \\
固定4段 & .551 & .665 & 4.00 \\
固定8段 & .623 & .713 & 8.00 \\
固定16段 & .660 & .735 & 16.00 \\
规划器 & .680 & .752 & 5.63 \\
规划器+回退 & .701 & .759 & 8.76 \\
\midrule
\multicolumn{5}{l}{\textit{GLM 5.3}} \\
固定4段 & .599 & .701 & 4.00 \\
固定8段 & .680 & .753 & 8.00 \\
固定16段 & .743 & .792 & 16.00 \\
规划器 & .726 & .774 & 5.69 \\
规划器+回退 & .753 & .789 & 9.19 \\
\midrule
\multicolumn{5}{l}{\textit{Qwen3-8B}} \\
固定4段 & .593 & .635 & 4.00 \\
固定8段 & .629 & .664 & 8.00 \\
固定16段 & .636 & .673 & 16.00 \\
规划器 & .637 & .657 & 5.12 \\
规划器+回退 & .649 & .663 & 6.13 \\
\bottomrule
\end{tabular}
\end{table}


表\ref{tab:findver}报告 FinDVer testmini 默认阅读器结果。DeepSeek Flash 上规划器加回退取得 0.701 准确率，平均 8.76 段；固定 16 段 0.660，固定 8 段 0.623。GLM 5.3 Flash 上该分支取得 0.753，平均 9.19 段；固定 16 段 0.743，固定 8 段 0.680。Qwen3-8B 上该分支取得 0.649，平均 6.13 段；固定 16 段 0.636，固定 8 段 0.629。分支在三个基座上均高于固定 16 段阅读器，平均少用 6.9 至 9.9 段。官方思维链阅读器在 DeepSeek Flash 和 GLM 5.3 Flash 上分别取得 0.817 和 0.819（金证据），十段检索版本均为 0.787；这些使用不同阅读器和证据接口，是上界参考，不是同条件基线。FinDVer 是迁移集，分支未在其上调参。

\begin{revbox}
\noindent\textbf{迁移的边界。}FinDVer 上 GLM 的宏 F1 为 .789，低于固定 16 段的 .792；Qwen 为 .663，低于 .673。准确率更高但宏 F1 回退，说明迁移收益不是全指标一致；正文只应把准确率作为控制分支的迁移证据，宏 F1 回退放入局限。Qwen3-8B 在 DocFinQA 上所有条件精确匹配低于 .11，主要原因是阅读器返回解释性文字而不是短答案；该基座不用于确定分支的主比较，只作为基座能力边界。
\end{revbox}

\section{全动作读出与公平基线}

\begin{revbox}
\noindent\textbf{核心结果。}本节把全动作损失感知读出应用到同一个后验和同一个损失，比较受限 gate 与全动作 gate；然后把同一个 gate 应用到所有具名基线，修正原稿的不公平比较。
\end{revbox}

\begin{table}[H]
\centering
\zihao{5}
\setlength{\tabcolsep}{3pt}
\caption{受限 gate 与全动作 gate 的决策损失比较（227 条 XBRL 核验断言；配对 95\% CI 为 10,000 次 claim 级 bootstrap）。}
\label{tab:fullaction}
\begin{tabular}{lrrrr}
\toprule
基座 & Fusion 受限 & Fusion 全动作 & RF 受限 & RF 全动作 \\
\midrule
DeepSeek & 3.604 & \cellcolor{yellow!25}2.863 & 3.454 & \cellcolor{yellow!25}2.881 \\
GLM 5.3 & 3.648 & \cellcolor{yellow!25}2.956 & 3.460 & \cellcolor{yellow!25}2.885 \\
Qwen3-8B & 3.577 & \cellcolor{yellow!25}2.837 & 3.454 & \cellcolor{yellow!25}2.881 \\
\bottomrule
\end{tabular}
\end{table}

\begin{revbox}
\noindent 全动作读出在三个基座上都显著降低 Fusion 的决策损失：DeepSeek 从 3.604 降到 2.863（差 $-0.740$，CI $[-1.031,-0.423]$），GLM 从 3.648 降到 2.956（差 $-0.690$，CI $[-1.000,-0.363]$），Qwen 从 3.577 降到 2.837（差 $-0.740$，CI $[-1.048,-0.423]$）。RF 上也分别从 3.454/3.460/3.454 降到 2.881/2.885/2.881。六个配对区间都排除零。这是本次修订中最稳健、最可复现的经验结果。
\end{revbox}

\begin{table}[H]
\centering
\zihao{5}
\setlength{\tabcolsep}{3pt}
\caption{公平 gate 下的基线比较（全动作读出；Fusion 为含 LLM 锚点信号的融合后验，RF 为结构化诊断）。}
\label{tab:fair}
\begin{tabular}{llrrr}
\toprule
基座 & 方法 & 决策损失 & 相对 Fusion & 95\% CI \\
\midrule
DeepSeek & 旧受限 Fusion & 3.604 & +0.741 & — \\
\rowcolor{yellow!18} & Fusion 全动作 & 2.863 & 0 & — \\
\rowcolor{yellow!18} & 最好 gated baseline（Self-RAG） & 3.048 & $-0.185$ & $[-0.581,0.211]$ \\
GLM 5.3 & 旧受限 Fusion & 3.648 & +0.692 & — \\
\rowcolor{yellow!18} & Fusion 全动作 & 2.956 & 0 & — \\
\rowcolor{yellow!18} & 最好 gated baseline（FinDVer） & 3.221 & $-0.265$ & $[-0.699,0.177]$ \\
Qwen3-8B & 旧受限 Fusion & 3.577 & +0.740 & — \\
\rowcolor{yellow!18} & Fusion 全动作 & 2.837 & 0 & — \\
\rowcolor{yellow!18} & 最好 gated baseline（Fin-Fact） & 4.846 & $-2.009$ & $[-2.687,-1.330]$ \\
\bottomrule
\end{tabular}
\end{table}

\begin{revbox}
\noindent 公平 gate 下，Fusion 的点估计在三个基座上都低于全部具名基线；Qwen 上显著，DeepSeek/GLM 上相对最强基线的差异不显著（CI 包含零）。因此论文应写“点估计更低、Qwen 与弱基线显著；最强基线在 DeepSeek/GLM 上不显著”，不能重复原稿“显著低于所有基线”的表述。旧受限 gate 的 Fusion 在 DeepSeek/GLM 上实际上高于最好 gated baseline；原稿的“最低损失”是“只有 FinPlanRAG 有 gate”造成的。
\end{revbox}

\section{演化版本：文本锚点与时间演化（DeepSeek）}

\begin{revbox}
\noindent\textbf{实验设置。}新 prompt 从截止时点 8-K 原文抽取带来源编号的锚点，公司名默认脱敏；时间演化只用揭示日早于 cutoff 的结果；演化先验 RF 在 baseline 特征上加入 peer/sector 历史、公司 surprise、日历和当前季度位置。融合权重只在 2019--2021 validation 上选择；2022--2025 test 只评估一次。
\end{revbox}

\begin{table}[H]
\centering
\zihao{5}
\setlength{\tabcolsep}{2.5pt}
\caption{DeepSeek 演化版系统与消融（227 条；Acc/Bal/Macro/Loss；Loss 为全动作 gate）。}
\label{tab:evolution}
\begin{tabular}{lrrrr}
\toprule
系统 & Acc & Bal & Macro & Loss \\
\midrule
旧论文 Fusion（受限 gate） & .524 & .471 & .468 & 3.604 \\
\rowcolor{yellow!18} 新锚点 + 演化 + 演化先验 RF（$w=.20$） & \textbf{.551} & \textbf{.489} & \textbf{.503} & \textbf{2.890} \\
\rowcolor{yellow!18} 演化先验 RF 单独（诊断） & \textbf{.555} & \textbf{.503} & \textbf{.508} & \textbf{2.828} \\
\rowcolor{yellow!18} 新锚点 + 演化 + baseline RF（$w=.05$） & .542 & .499 & .497 & 2.855 \\
\rowcolor{yellow!18} 新锚点，静态先验/可靠性（$w=.05$） & .546 & .505 & .503 & 2.908 \\
\rowcolor{yellow!18} 旧锚点 + 演化（$w=.30$） & .520 & .439 & .438 & 2.793 \\
\rowcolor{yellow!18} 旧锚点，静态（$w=.20$） & .507 & .450 & .450 & 2.899 \\
\bottomrule
\end{tabular}
\end{table}

\begin{revbox}
\noindent 最终系统的 Acc/Bal/Macro 从旧系统的 .524/.471/.468 提高到 .551/.489/.503，Loss 从 3.604 降到 2.890。配对损失差为 $-0.714$，CI $[-1.057,-0.352]$，显著；但相对演化先验 RF 诊断的损失差为 $+0.062$，CI $[-0.141,0.273]$，不显著。准确率相对旧系统 +2.7 点，McNemar $p\approx0.69$，只能作为点估计报告。旧锚点 + 演化的 Loss 最低（2.793），但 Bal/Macro 只有 .439/.438；新锚点系统是更好的准确率/损失折中，也是更可辩护的（文本来源、脱敏）配置。
\end{revbox}

\begin{table}[H]
\centering
\zihao{5}
\setlength{\tabcolsep}{3pt}
\caption{演化先验特征（无 LLM）的 validation/test 结果。Validation 为 185 条；test 为 227 条；相同 3 seeds/300 trees。}
\label{tab:prior}
\begin{tabular}{llrrrr}
\toprule
Split & 特征集 & Acc & Bal & Macro & Loss \\
\midrule
Validation & baseline & .557 & .556 & .536 & 3.784 \\
\rowcolor{yellow!18} Validation & all evolved, no artifact & \textbf{.605} & \textbf{.572} & \textbf{.566} & \textbf{3.740} \\
Test & baseline RF（受限 / 全动作） & .520 & .476 & .473 & 3.463 / 2.908 \\
\rowcolor{yellow!18} Test & evolved-prior RF（受限 / 全动作） & \textbf{.555} & \textbf{.503} & \textbf{.508} & \textbf{3.278 / 2.828} \\
\bottomrule
\end{tabular}
\end{table}

\begin{revbox}
\noindent 演化先验把测试准确率从 .520 提到 .555（+3.5 点，McNemar $p\approx0.40$），平衡准确率 .476→.503，宏 F1 .473→.508；这些是点估计，不是显著性主张。Peer-only 特征在 validation 上反而使准确率降到 .524，说明增益来自 peer、公司 surprise、日历和当前季度位置的组合，不能归因于单一 peer 先验。\texttt{current\_eq\_low} 是数据构造痕迹，已从最终特征集中排除。
\end{revbox}

\section{消融与可信集}

\begin{revbox}
\noindent\textbf{全动作 vs 受限 gate。}表\ref{tab:fullaction}已经给出：Fusion/RF 在三个基座上都是全动作更好，且 CI 排除零。\textbf{演化 vs 静态。}新锚点 static 2.908→evolved 2.855；旧锚点 static 2.899→evolved 2.793；方向一致，但幅度小，227 条上不显著。\textbf{新 vs 旧锚点。}新锚点的 Acc/Bal/Macro 更好（.551/.489/.503 对 .520/.439/.438），旧锚点的 Loss 更低；两者是不同权衡。\textbf{演化先验 vs baseline RF。}演化先验在 Acc/Bal/Macro 上更好，Loss 在全动作下从 2.908 降到 2.828；相对旧受限系统显著，相对新锚点融合不显著。
\end{revbox}

\begin{table}[H]
\centering
\zihao{5}
\setlength{\tabcolsep}{3pt}
\caption{动作集与可信集（DeepSeek 227 条）。}
\label{tab:sets}
\begin{tabular}{lrrr}
\toprule
对象 & 平均集合大小 & 覆盖率 & 说明 \\
\midrule
$1-\alpha$ 可信集（Fusion） & 2.83 & .991 & 审计/覆盖口径 \\
\rowcolor{yellow!18} 全动作裁决集（Fusion） & 1.90 & .868 & 损失最优动作 \\
$1-\alpha$ 可信集（RF） & 2.76 & .987 & 审计/覆盖口径 \\
\rowcolor{yellow!18} 全动作裁决集（RF） & 1.89 & .863 & 损失最优动作 \\
\bottomrule
\end{tabular}
\end{table}

\begin{revbox}
\noindent 全动作裁决集的平均大小约 1.9，覆盖率约 .86，低于 $1-\alpha$ 集的 .99；这是给定损失（漏报 8、宽度 2）下的最优权衡，不是覆盖率缺陷。论文应同时报告 $1-\alpha$ 可信集（覆盖/审计）和损失最优裁决集（决策），不要把两者混为一个“可信集”。DeepSeek 上全动作动作全部为 adjudicate，没有 answer/abstain；这是当前损失参数下的动作分布，换损失参数会改变动作组合。
\end{revbox}

\section{误差分析与讨论}

Router 的错误集中在语义边界。DeepSeek Flash 和 GLM 5.3 Flash 把“未来支付金额”“预期股利”等确定问题判为前瞻性断言；这些问题的答案在文件中已经给出，但文本包含未来时态。Qwen3-8B 的错误更多，说明小模型的语义边界判断不稳定。错路由的置信度在 DeepSeek 上是 .9/.85/.85，属于高置信度错误；原稿“低置信度时人工复核”的补救机制抓不到这些错误。后续需要置信度校准或轻量检索信号，而不是简单的阈值门。

AnchorEvolve 的主要误差来自类别偏置和 LLM 抽取器的主观可靠性。纯 LLM 锚点后验在旧 prompt 下准确率 .445，新 prompt 下 .502，仍低于结构化 RF；锚点层不是点预测的主因。融合后系统的点估计与 RF 接近，损失相对旧系统显著下降，但相对 RF 诊断的增量不显著。文本锚点的价值在于把 8-K 原文中的管理层评论、分部/需求信号变成可审计的 L2 状态，而不是替代结构化先验。

\begin{revbox}
\noindent\textbf{外部效度与污染。}新文本锚点目前只跑了 DeepSeek；GLM/Qwen 只验证了全动作读出（旧锚点），因为需要对应的 API key/本地模型。收入指引标签来自 XBRL 自动核验，不是人工金标准；227 条测试中 33 条为 below，类别不平衡使平衡准确率/宏 F1 的点估计不稳定。公司名已脱敏，但 LLM 预训练记忆仍可能污染 2022--2025 的结果；需要用 cutoff 早于 2022 的模型或合成公司做对照。46 条纯文本标签的敏感性结果、FinDVer 的宏 F1 回退、DocFinQA 的 300/7,437 样本比例和 planner-16 支配关系都应在局限/附录中报告。
\end{revbox}

% ==================== 第7章（重构版） ====================
\chapter{讨论与结论}
\label{ch:conclusion}

\begin{revbox}
\noindent\textbf{本章修订说明。}本章按修订后的证据层级重写讨论与结论：router 是接口，确定分支是控制，AnchorEvolve 是可演化的证据状态，全动作损失感知读出是核心经验贡献。原稿“LLM 锚点不可替代”“融合后验取得最高准确率”等表述不再保留；准确率相对 RF 的 +2.7 点只作为点估计，显著结果是相对旧受限 gate 的损失下降。
\end{revbox}

\section{讨论}

路由实验回答第一个研究问题。DeepSeek Flash 和 GLM 5.3 Flash 分别达到 0.995 和 0.990 的路由准确率，混合效用为 0.403 和 0.419，后者恢复 oracle 的 99.3\%。错路由主要发生在包含未来时态、但答案已经在文件中给出的确定问题；错路由的置信度较高（DeepSeek 上为 .9/.85/.85），说明简单的低置信度人工复核抓不到这些错误。Qwen3-8B 的路由准确率降至 0.867，混合效用 0.263，低于 Always-Anchor 的 0.298。这个对比说明契约识别不是简单的时态匹配：较强基座能够利用查询与截止时点的关系，较小基座则容易把语义边界压缩为表面模式。

\begin{revbox}
\noindent\textbf{Router 的定位。}本文不主张 router 结构创新：它是一个零样本 LLM 接口，用来把契约选择显式化。H1 检验的是“契约选择是否可学”，不是“router 分类器是否优于 Adaptive-RAG/Self-RAG 等已有路由”。混合效用的提升来自分支契约的改进（527 条已验证子集上 .389→.437），而不是 router 本身的重新设计。若论文把 router 当第一贡献，审稿人只会看到增量；把它写成接口，才能把创新让给契约形式化、演化状态和读出层。
\end{revbox}

DocFinQA 和 FinDVer 回答第二个研究问题，但它们是\textbf{控制分支}的结果。DeepSeek 上，义务检索分支取得 0.457 精确匹配、平均 10.69 段，固定 16 段阅读器取得 0.390、平均 16 段；GLM 上分别为 0.497 对 0.480，平均 10.31 段对 16 段。FinDVer 上三个基座分别为 0.701、0.753 和 0.649，均高于固定 16 段的 0.660、0.743 和 0.636。四个具名 RAG 适配在 DocFinQA 上均低于义务检索分支，分支的答案覆盖率 0.770 和 0.810 也高于固定预算。

\begin{revbox}
\noindent\textbf{控制分支的边界。}确定分支不是本文的新检索算法；主实验使用确定性的词项义务覆盖，类型化 metadata 过滤是设计目标而非已实现性质。成本口径也要诚实：DeepSeek 上最终分支相对固定 16 段只省 2.3\% 的总 token，且 planner B=16 在准确率和 token 上都支配最终分支；GLM 上最终分支省 29\% token。FinDVer 上 GLM/Qwen 的宏 F1 低于固定 16 段。因此 H2 只支持“控制分支足以让 router 的另一半可信”，不支持“新检索算法一致占优”。
\end{revbox}

收入指引实验回答第三个研究问题，也是本次修订的核心。旧系统（受限 gate）的 Fusion 决策损失为 3.604；全动作损失感知读出在同一后验和同一损失下把它降到 2.863（DeepSeek）、2.956（GLM）和 2.837（Qwen），三个配对区间都排除零。RF 上也从 3.454/3.460/3.454 降到 2.881/2.885/2.881。该读出是严格改进：原 gate 只在点答案与最高密度 $1-\alpha$ 可信集之间选择，新 gate 枚举全部点标签、7 个非空裁决集和弃权。

\begin{revbox}
\noindent\textbf{公平基线。}把所有具名基线放到同一个全动作 gate 下，Fusion 的点估计低于全部基线；Qwen 上显著，DeepSeek/GLM 上相对最强基线不显著（DeepSeek 对 Self-RAG 差 $-0.185$，CI $[-0.581,0.211]$；GLM 对 FinDVer 差 $-0.265$，CI $[-0.699,0.177]$）。原稿“显著低于所有基线”只成立于“baseline 只能点答案、FinPlanRAG 可以用集合/弃权”的不对等设定；修订版不再使用这个比较。
\end{revbox}

文本锚点与时间演化给出\textbf{增量但不是主因}的结果。新 prompt 把 LLM 单独准确率从 .445 提到 .502；新锚点 + 演化 + 演化先验 RF 的最终系统达到 Acc .551、Bal .489、Macro .503、Loss 2.890，相对旧系统损失差 $-0.714$，CI $[-1.057,-0.352]$，显著；相对演化先验 RF 诊断（.555/.503/.508/2.828）的损失差为 $+0.062$，CI $[-0.141,0.273]$，不显著。准确率相对旧系统 +2.7 点，但 McNemar $p\approx0.69$，只能作为点估计。

\begin{revbox}
\noindent\textbf{演化消融。}新锚点 static 2.908→evolved 2.855；旧锚点 static 2.899→evolved 2.793；方向一致但幅度小，227 条上不显著。演化先验 RF 把测试准确率从 .520 提到 .555（+3.5 点，McNemar $p\approx0.40$），Bal/Macro 分别 .476→.503、.473→.508；这些也是点估计。结论应写成“时间演化和演化先验提供了一致的点估计改善，显著增益来自全动作读出”，而不是“演化是主因”。
\end{revbox}

\section{局限与未来研究}

当前证据边界主要来自五点。第一，新文本锚点只跑了 DeepSeek；GLM/Qwen 只验证了全动作读出（旧锚点），需要对应 API key/本地模型才能完成三基座的文本锚点实验。第二，收入指引标签来自 XBRL 自动核验，不是人工金标准；227 条测试中 33 条 below，类别不平衡使平衡准确率/宏 F1 的点估计不稳定；46 条纯文本标签的敏感性结果仍需报告。第三，公司名已脱敏，但 LLM 预训练记忆仍可能污染 2022--2025 的结果；需要用 cutoff 早于 2022 的模型或合成公司做对照。第四，“可审计”是设计属性，尚未做人工审计实验；锚点来源、资格理由和信念变化目前只是可检查的接口，不是已测量的用户收益。第五，全动作裁决集覆盖率约 .86，低于 $1-\alpha$ 集的 .99；这是给定损失下的最优权衡，但不同机构换损失参数后动作分布会改变，需要敏感性分析。

\begin{revbox}
\noindent\textbf{控制分支与实现边界。}确定分支的主结果仍使用 regex 词项覆盖和固定 chunk 时间戳，不是第 4 章的类型化 planner；DocFinQA 只跑 300/7,437 样本；planner B=16 在 DeepSeek 上支配最终分支；FinDVer 的 GLM/Qwen 宏 F1 回退。这些是控制分支的边界，不应被包装成主贡献。类型化 planner、全测试集、dense/hybrid 检索和预算对齐是确定分支的后续工作，不应挤占 AnchorEvolve/readout 的主线。
\end{revbox}

后续研究将围绕五个方向展开。第一，做 LLM 污染检验：公司名脱敏之外的 cutoff 早于 2022 模型、合成公司、或“是否知道结果”的对照。第二，把时间演化扩展到多轮/滚动 origin，并在更多前瞻性事件（生产里程碑、产能投放、监管审批、合同履约）上评估，而不是只做收入指引三桶。第三，补人工审计实验：有/无锚点 provenance 时，审计错误发现率/时间/信心是否改善；如果“可审计”要继续作为贡献，这是必须的。第四，补简单 router 基线（关键词/时点规则、TF-IDF/逻辑回归、统一灵活输出接口的 LLM），把 H1 从“LLM 能分”收紧为“契约路由比简单规则/统一接口更好”。第五，做损失参数、$\alpha$、$w$、$\rho$、$p_e$、硬锚点阈值的敏感性网格，并报告 coverage-size 曲线与 conformal/阈值 gate 对照。

\section{研究结论}

本文把金融混合查询的回答形式选择表述为查询契约的查询级决策，并以义务检索和 AnchorEvolve 分别实现闭合答案契约和前瞻性判断契约。Router 是这一契约选择器的简单接口，不是本文的主要算法创新。确定契约由义务检索控制分支实现；前瞻契约由 AnchorEvolve 实现：从截止时点 8-K 原文抽取带来源的 L2 锚点，经五类资格判定、时间演化的类型可靠性/先验和信念更新得到可审计证据状态，再由全动作损失感知读出转换为点答案、裁决集或弃权。

\begin{revbox}
\noindent 三项贡献的证据层级是：\textbf{契约级形式化}是概念贡献，router 只是接口；\textbf{可演化的锚点证据状态}是状态贡献，其价值在审计、时点正确和作为决策输入，而不是单独点预测；\textbf{全动作损失感知读出}是核心经验贡献，在 DeepSeek、GLM、Qwen 三个基座上显著降低决策损失（3.604→2.863、3.648→2.956、3.577→2.837），并把基线比较放到同一个 gate 下。DeepSeek 演化版最终系统达到 .551/.489/.503/2.890，相对旧系统显著降损（$-0.714$，CI $[-1.057,-0.352]$）；演化先验 RF 诊断达到 .555/.503/.508/2.828。准确率相对 RF 的 +2.7 点不显著，论文不把它写成显著优势。
\end{revbox}

从金融实践看，框架把答案、证据和风险放在同一条审计链上。分析师对确定问题获得带证据编号的闭合答案，对前瞻性断言获得点判断、损失最优裁决集或弃权，并可以查看锚点来源、资格理由和信念变化。机构可以按自身风险偏好替换损失参数和校准后端，而不改变契约选择与证据状态的结构。生产里程碑、产能投放、监管审批和合同履约等任务与收入指引具有相似的时点结构，本文结果为这些场景提供可复用的接口。

本文的结论有明确边界。收入指引实验采用三桶结果空间；文本锚点目前只在 DeepSeek 上完成；标签来自 XBRL 自动核验；227 条测试使准确率/平衡准确率的点估计不稳定；全动作裁决集覆盖率低于 $1-\alpha$ 集；可审计性尚未做人工实验。后续工作将围绕污染检验、多轮演化、人工审计、简单 router 基线和损失参数敏感性展开。由于契约选择、锚点状态和读出彼此分离，这些改进可以独立评估，而不需要改变 FinPlanRAG 的整体接口。


\clearpage
\phantomsection
\addcontentsline{toc}{chapter}{参考文献}
{\zihao{5}\setlength{\baselineskip}{16pt}\begin{thebibliography}{33}
\setlength{\bibsep}{0pt}
\setlength{\itemsep}{0pt}
\setlength{\parsep}{0pt}
\setlength{\bibhang}{2em}

\bibitem[Angelopoulos and Bates(2021)]{angelopoulos2021gentle}
Angelopoulos, A. N., and Bates, S., 2021, ``A Gentle Introduction to Conformal Prediction and Distribution-Free Uncertainty Quantification'', arXiv:2107.07511.

\bibitem[Asai et al.(2024)]{asai2024selfrag}
Asai, A., Wu, Z., Wang, Y., Sil, A., and Hajishirzi, H., 2024, ``Self-RAG: Learning to Retrieve, Generate, and Critique through Self-Reflection'', in \emph{International Conference on Learning Representations}.

\bibitem[Brodersen et al.(2010)]{brodersen2010balanced}
Brodersen, K. H., Ong, C. S., Stephan, K. E., and Buhmann, J. M., 2010, ``The Balanced Accuracy and Its Posterior Distribution'', in \emph{Proceedings of the 20th International Conference on Pattern Recognition}, pp. 3121--3124.

\bibitem[Chen et al.(2026)]{chen2026hcrag}
Chen, S., Tan, H., Li, Y., and Liang, J., 2026, ``HC-RAG: Evidence-Centric Retrieval-Augmented Generation over Heterogeneous Financial Filings'', arXiv:2608.12335.

\bibitem[Chen et al.(2021)]{chen2021finqa}
Chen, Z., Chen, W., Smiley, C., Shah, S., Borova, I., Langdon, D., Moussa, R., Beane, M., Huang, T.-H., Routledge, B., and Wang, W. Y., 2021, ``FinQA: A Dataset of Numerical Reasoning over Financial Data'', in \emph{Proceedings of the 2021 Conference on Empirical Methods in Natural Language Processing}, pp. 3697--3711.

\bibitem[Chi et al.(2026)]{chi2026finsagent}
Chi, J., Tai, Z., Wu, H., Kwok, T. S. T., He, H., Liao, Z., Xiao, B., Jiang, C., Lei, J., Huang, J., Lu, P., Li, M., Ma, L., Wu, Y., Lyu, S., Tian, J., Li, Y., Ma, Y., Guan, S., Hu, D., Cui, Y., Zhou, L., Ding, L., and Wang, X., 2026, ``FinSAgent: Corpus-Aligned Multi-Agent RAG Framework for Evidence-Grounded SEC Filing Question Answering'', arXiv:2607.18102.

\bibitem[Chow(1970)]{chow1970optimum}
Chow, C. K., 1970, ``On Optimum Recognition Error and Reject Tradeoff'', \emph{IEEE Transactions on Information Theory}, Vol. 16, No. 1, pp. 41--46.

\bibitem[Elkan(2001)]{elkan2001cost}
Elkan, C., 2001, ``The Foundations of Cost-Sensitive Learning'', in \emph{Proceedings of the 17th International Joint Conference on Artificial Intelligence}, pp. 973--978.

\bibitem[Geifman and El-Yaniv(2017)]{geifman2017selective}
Geifman, Y., and El-Yaniv, R., 2017, ``Selective Classification for Deep Neural Networks'', in \emph{Advances in Neural Information Processing Systems}, Vol. 30.

\bibitem[Gibbs and Cand\`es(2021)]{gibbs2021adaptive}
Gibbs, I., and Cand\`es, E. J., 2021, ``Adaptive Conformal Inference Under Distribution Shift'', in \emph{Advances in Neural Information Processing Systems}, Vol. 34, pp. 1660--1672.

\bibitem[Guo et al.(2017)]{guo2017calibration}
Guo, C., Pleiss, G., Sun, Y., and Weinberger, K. Q., 2017, ``On Calibration of Modern Neural Networks'', in \emph{Proceedings of the 34th International Conference on Machine Learning}, Vol. 70, pp. 1321--1330.

\bibitem[Hirst et al.(2008)]{hirst2008management}
Hirst, D. E., Koonce, L., and Venkataraman, S., 2008, ``Management Earnings Forecasts: A Review and Framework'', \emph{Accounting Horizons}, Vol. 22, No. 3, pp. 315--338.

\bibitem[Islam et al.(2023)]{islam2023financebench}
Islam, P., Kannappan, A., Kiela, D., Qian, R., Scherrer, N., and Vidgen, B., 2023, ``FinanceBench: A New Benchmark for Financial Question Answering'', arXiv:2311.11944.

\bibitem[Jeong et al.(2024)]{jeong2024adaptiverag}
Jeong, S., Baek, J., Cho, S., Hwang, S. J., and Park, J. C., 2024, ``Adaptive-RAG: Learning to Adapt Retrieval-Augmented Large Language Models through Question Complexity'', in \emph{Proceedings of the 2024 Conference of the North American Chapter of the Association for Computational Linguistics}, pp. 7036--7050.

\bibitem[Jiang et al.(2026)]{jiang2026finrate}
Jiang, Y., Chen, J., Makri, E., Chen, J., Li, P., Maatouk, A., Tassiulas, L., Brenner, E., Xiang, B., and Ying, R., 2026, ``Fin-RATE: A Real-world Financial Analytics and Tracking Evaluation Benchmark for LLMs on SEC Filings'', in \emph{Proceedings of the 32nd ACM SIGKDD Conference on Knowledge Discovery and Data Mining}, pp. 9183--9194.

\bibitem[Jiang et al.(2023)]{jiang2023flare}
Jiang, Z., Xu, F., Gao, L., Sun, Z., Liu, Q., Dwivedi-Yu, J., Yang, Y., Callan, J., and Neubig, G., 2023, ``Active Retrieval Augmented Generation'', in \emph{Proceedings of the 2023 Conference on Empirical Methods in Natural Language Processing}, pp. 7969--7992.

\bibitem[Lewis et al.(2020)]{lewis2020rag}
Lewis, P., Perez, E., Piktus, A., Petroni, F., Karpukhin, V., Goyal, N., Kuttler, H., Lewis, M., Yih, W.-T., Rockt\"aschel, T., Riedel, S., and Kiela, D., 2020, ``Retrieval-Augmented Generation for Knowledge-Intensive NLP Tasks'', in \emph{Advances in Neural Information Processing Systems}, Vol. 33, pp. 9459--9474.

\bibitem[L\'opez de Prado(2018)]{lopezdeprado2018advances}
L\'opez de Prado, M., 2018, \emph{Advances in Financial Machine Learning}, Wiley.

\bibitem[Manski(2003)]{manski2003partial}
Manski, C. F., 2003, \emph{Partial Identification of Probability Distributions}, Springer.

\bibitem[Mansouri et al.(2026)]{mansouri2026finrank}
Mansouri, S., Saad, D., Wahrenburg, M., Weissel, M., and Woebbeking, F., 2026, ``FinRank: An Evidence-Grounded Benchmark for Financial Question Answering and Retrieval over SEC Filings'', arXiv:2608.07400.

\bibitem[Qwen Team(2025)]{qwen2025qwen3}
Qwen Team, 2025, ``Qwen3 Technical Report'', arXiv:2505.09388.

\bibitem[Rangapur et al.(2025)]{rangapur2025finfact}
Rangapur, A., Wang, H., Jian, L., and Shu, K., 2025, ``Fin-Fact: A Benchmark Dataset for Multimodal Financial Fact-Checking and Explanation Generation'', in \emph{Companion Proceedings of the ACM Web Conference 2025}, pp. 785--788.

\bibitem[Reddy et al.(2024)]{reddy2024docfinqa}
Reddy, V., Koncel-Kedziorski, R., Lai, V. D., Krumdick, M., Lovering, C., and Tanner, C., 2024, ``DocFinQA: A Long-Context Financial Reasoning Dataset'', in \emph{Proceedings of the 62nd Annual Meeting of the Association for Computational Linguistics (Volume 2: Short Papers)}, pp. 445--458.

\bibitem[Romano et al.(2020)]{romano2020classification}
Romano, Y., Sesia, M., and Cand\`es, E. J., 2020, ``Classification with Valid and Adaptive Coverage'', in \emph{Advances in Neural Information Processing Systems}, Vol. 33, pp. 3581--3591.

\bibitem[Shinn et al.(2023)]{shinn2023reflexion}
Shinn, N., Cassano, F., Berman, A., Gopinath, A., Narasimhan, K., and Yao, S., 2023, ``Reflexion: Language Agents with Verbal Reinforcement Learning'', in \emph{Advances in Neural Information Processing Systems}, Vol. 36.

\bibitem[Sokolova and Lapalme(2009)]{sokolova2009performance}
Sokolova, M., and Lapalme, G., 2009, ``A Systematic Analysis of Performance Measures for Classification Tasks'', \emph{Information Processing and Management}, Vol. 45, No. 4, pp. 427--437.

\bibitem[Vovk et al.(2005)]{vovk2005algorithmic}
Vovk, V., Gammerman, A., and Shafer, G., 2005, \emph{Algorithmic Learning in a Random World}, Springer.

\bibitem[Xie et al.(2024)]{xie2024finben}
Xie, Q., Han, W., Chen, Z., Xiang, R., Zhang, X., He, Y., Xiao, M., Li, D., Dai, Y., Feng, D., Xu, Y., Kang, H., Kuang, Z., Yuan, C., Yang, K., Luo, Z., Zhang, T., Liu, Z., Xiong, G., Deng, Z., Jiang, Y., Yao, Z., Li, H., Yu, Y., Hu, G., Huang, J., Liu, X.-Y., Lopez-Lira, A., Wang, B., Lai, Y., Wang, H., Peng, M., Ananiadou, S., and Huang, J., 2024, ``FinBen: A Holistic Financial Benchmark for Large Language Models'', in \emph{Advances in Neural Information Processing Systems}, Vol. 37.

\bibitem[Yan et al.(2024)]{yan2024crag}
Yan, S.-Q., Gu, J.-C., Zhu, Y., and Ling, Z.-H., 2024, ``Corrective Retrieval Augmented Generation'', arXiv:2401.15884.

\bibitem[Yao et al.(2023)]{yao2023react}
Yao, S., Zhao, J., Yu, D., Du, N., Shafran, I., Narasimhan, K., and Cao, Y., 2023, ``ReAct: Synergizing Reasoning and Acting in Language Models'', in \emph{International Conference on Learning Representations}.

\bibitem[Zhao et al.(2024)]{zhao2024findver}
Zhao, Y., Long, Y., Jiang, T., Wang, C., Chen, W., Liu, H., Tang, X., Zhang, Y., Zhao, C., and Cohan, A., 2024, ``FinDVer: Explainable Claim Verification over Long and Hybrid-Content Financial Documents'', in \emph{Proceedings of the 2024 Conference on Empirical Methods in Natural Language Processing}, pp. 14739--14752.

\bibitem[Zhou et al.(2026)]{zhou2026fincards}
Zhou, Y., Zhang, F., Chen, Y., Zhang, H., Nakov, P., and Xie, Z., 2026, ``FinCARDS: Card-Based Analyst Reranking for Financial Document Question Answering'', in \emph{Findings of the Association for Computational Linguistics: ACL 2026}, pp. 24836--24852.

\bibitem[Zhu et al.(2021)]{zhu2021tatqa}
Zhu, F., Lei, W., Huang, Y., Wang, C., Zhang, S., Lv, J., Feng, F., and Chua, T.-S., 2021, ``TAT-QA: A Question Answering Benchmark on a Hybrid of Tabular and Textual Content in Finance'', in \emph{Proceedings of the 59th Annual Meeting of the Association for Computational Linguistics}, pp. 3277--3287.

\end{thebibliography}
}

% ==================== 附录 ====================
\appendix
\chapter{完整指标表}
\label{app:fullmetrics}

表\ref{tab:vague-ds}至表\ref{tab:vague-qwen}报告三个基座上的完整十一项指标。主文表\ref{tab:fullaction}和表\ref{tab:evolution}报告全动作读出与演化版结果；本附录表 A.1--A.3 保留旧受限 gate 的完整十一项指标，用于历史对照。完整指标中的AUROC和AUPRC衡量排序质量，ECE和Brier衡量概率校准，Sel@50和AURC衡量选择性预测，决策损失衡量非对称风险。
ef{tab:vague-ds}至表
ef{tab:vague-qwen}报告三个基座上的完整十一项指标。主文表
ef{tab:fullaction}和表
ef{tab:evolution}报告全动作读出与演化版结果；本附录表 A.1--A.3 保留旧受限 gate 的完整十一项指标，用于历史对照。完整指标中的AUROC和AUPRC衡量排序质量，ECE和Brier衡量概率校准，Sel@50和AURC衡量选择性预测，决策损失衡量非对称风险。

\begin{table}[htbp]
\centering
\zihao{5}
\setlength{\tabcolsep}{2pt}
\caption{DeepSeek基座上的模糊断言结果（227条XBRL核验断言）。Loss越低越好；Sel@50为50\%覆盖率下的选择性准确率，AURC越低越好。}
\label{tab:vague-ds}
\resizebox{\linewidth}{!}{%
\begin{tabular}{lrrrrrrrrrrr}
\toprule
方法 & Acc & Bal & Macro & MCC & AUROC & AUPRC & ECE & Brier & Loss & Sel@50 & AURC \\
\midrule
FinDVer-style & \textbf{.551} & .461 & .464 & \textbf{.233} & .674 & .494 & .092 & .580 & 4.678 & .675 & .337 \\
Fin-Fact-style & .537 & .446 & .450 & .201 & .650 & .475 & .075 & .590 & 4.705 & .675 & .338 \\
Self-RAG-style & .546 & .451 & .456 & .216 & .670 & .491 & .102 & .580 & 4.564 & \textbf{.684} & \textbf{.333} \\
CRAG-style & .529 & .423 & .418 & .177 & .643 & .473 & .097 & .593 & 4.802 & .658 & .349 \\
FLARE-style & .511 & .417 & .411 & .164 & .634 & .467 & .130 & .610 & 4.881 & .596 & .362 \\
Adaptive-RAG-style & .493 & .406 & .400 & .143 & .584 & .425 & .147 & .634 & 5.093 & .579 & .415 \\
AnchorEvolve-LLM & .445 & .307 & .242 & -0.075 & .465 & .312 & .129 & .666 & 4.035 & .509 & .529 \\
FinPlanRAG（AnchorEvolve-Fusion, $w=0.1$） & .524 & .471 & .468 & .201 & .703 & .503 & \textbf{.033} & \textbf{.564} & 3.604 & .632 & .383 \\
共享校准RF（非LLM诊断） & .533 & \textbf{.491} & \textbf{.486} & .229 & \textbf{.706} & \textbf{.505} & .033 & .566 & \textbf{3.454} & .649 & .383 \\
\bottomrule
\end{tabular}
}
\end{table}


\begin{table}[htbp]
\centering
\zihao{5}
\setlength{\tabcolsep}{2pt}
\caption{GLM 5.3基座上的模糊断言结果（227条XBRL核验断言）。Loss越低越好；Sel@50为50\%覆盖率下的选择性准确率，AURC越低越好。}
\label{tab:vague-glm}
\resizebox{\linewidth}{!}{%
\begin{tabular}{lrrrrrrrrrrr}
\toprule
方法 & Acc & Bal & Macro & MCC & AUROC & AUPRC & ECE & Brier & Loss & Sel@50 & AURC \\
\midrule
FinDVer-style & .498 & .439 & .442 & .178 & .642 & .439 & .103 & .627 & 5.233 & .605 & .418 \\
Fin-Fact-style & .454 & .398 & .402 & .089 & .603 & .430 & .178 & .670 & 5.339 & .579 & .427 \\
Self-RAG-style & .454 & .406 & .404 & .113 & .617 & .434 & .164 & .637 & 5.163 & .588 & .411 \\
CRAG-style & .515 & .431 & .435 & .173 & .607 & .418 & .148 & .657 & 4.846 & .570 & .441 \\
FLARE-style & .471 & .388 & .387 & .110 & .583 & .389 & .192 & .705 & 5.286 & .526 & .518 \\
Adaptive-RAG-style & .427 & .372 & .376 & .059 & .580 & .382 & .227 & .738 & 5.630 & .482 & .542 \\
AnchorEvolve-LLM & .454 & .329 & .295 & -0.004 & .498 & .339 & .145 & .671 & 3.903 & .430 & .552 \\
FinPlanRAG（AnchorEvolve-Fusion, $w=0.1$） & \textbf{.537} & .489 & \textbf{.487} & .225 & .704 & .504 & \textbf{.032} & \textbf{.562} & 3.648 & .640 & \textbf{.376} \\
共享校准RF（非LLM诊断） & .533 & \textbf{.491} & .486 & \textbf{.229} & \textbf{.706} & \textbf{.505} & .033 & .566 & \textbf{3.454} & \textbf{.649} & .383 \\
\bottomrule
\end{tabular}
}
\end{table}


\begin{table}[htbp]
\centering
\zihao{5}
\setlength{\tabcolsep}{2pt}
\caption{Qwen3-8B基座上的模糊断言结果（227条XBRL核验断言）。Loss越低越好；Sel@50为50\%覆盖率下的选择性准确率，AURC越低越好。}
\label{tab:vague-qwen}
\resizebox{\linewidth}{!}{%
\begin{tabular}{lrrrrrrrrrrr}
\toprule
方法 & Acc & Bal & Macro & MCC & AUROC & AUPRC & ECE & Brier & Loss & Sel@50 & AURC \\
\midrule
FinDVer-style & .445 & .422 & .406 & .133 & .564 & .375 & .555 & 1.110 & 5.524 & .456 & .535 \\
Fin-Fact-style & .511 & .477 & .464 & \textbf{.237} & .609 & .404 & .489 & .978 & 4.846 & .491 & .467 \\
Self-RAG-style & .489 & .447 & .437 & .196 & .583 & .394 & .489 & .990 & 5.093 & .482 & .506 \\
CRAG-style & .454 & .418 & .399 & .152 & .566 & .371 & .547 & 1.093 & 5.366 & .456 & .534 \\
FLARE-style & .326 & .298 & .212 & -0.107 & .475 & .322 & .672 & 1.340 & 6.740 & .316 & .664 \\
Adaptive-RAG-style & .344 & .324 & .290 & -0.027 & .492 & .330 & .654 & 1.308 & 6.485 & .298 & .663 \\
AnchorEvolve-LLM & .427 & .353 & .355 & .027 & .562 & .399 & .163 & .647 & 3.586 & .561 & .435 \\
FinPlanRAG（AnchorEvolve-Fusion, $w=0.1$） & \textbf{.537} & .486 & .485 & .227 & \textbf{.707} & \textbf{.509} & \textbf{.027} & \textbf{.562} & 3.577 & .632 & \textbf{.373} \\
共享校准RF（非LLM诊断） & .533 & \textbf{.491} & \textbf{.486} & .229 & .706 & .505 & .033 & .566 & \textbf{3.454} & \textbf{.649} & .383 \\
\bottomrule
\end{tabular}
}
\end{table}


\chapter{Qwen3-8B的DocFinQA结果}
\label{app:qwen}

Qwen3-8B在DocFinQA上返回解释性文字而非短答案，精确匹配接近地板。表\ref{tab:qwen-doc}报告该基座的全部条件；它不用于义务检索分支的主比较，只作为基座能力边界。该结果说明，短答案抽取任务对阅读器的格式遵循能力有额外要求；当阅读器返回解释性文字时，精确匹配和词级F1会同时下降，但答案覆盖率仍可能较高。

\begin{table}[htbp]
\centering
\zihao{5}
\setlength{\tabcolsep}{4pt}
\caption{Qwen3-8B的DocFinQA结果（300题固定样本）。}
\label{tab:qwen-doc}
\begin{tabular}{lrrrr}
\toprule
条件 & 准确率 & 宏F1 & 平均段落 & 覆盖率 \\
\midrule
固定4段 & .100 & .048 & 4.00 & .760 \\
固定8段 & .110 & .040 & 8.00 & .833 \\
固定16段 & .100 & .046 & 16.00 & .853 \\
规划器B=8 & .093 & .034 & 3.67 & .883 \\
规划器+回退 & .093 & .034 & 5.85 & .883 \\
\bottomrule
\end{tabular}
\end{table}

\chapter{命题与完整证明}
\label{app:proofs}

本章给出正文使用的四个命题及完整证明。证明只使用有限结果空间、有限动作集和线性期望，不依赖具体数据或模型。

\section{裁决门的贝叶斯最优性}

\begin{proposition}
给定信念$b\in\Delta(\mathcal H)$、动作集合$\mathcal A=\{\text{answer},\text{adjudicate},\text{abstain}\}$和损失函数$L:\mathcal A\times\mathcal H\to\mathbb R$，选择
\[
\pi^\star=\arg\min_{\pi\in\mathcal A}\mathbb E_{y\sim b}[L(\pi,y)]
\]
是贝叶斯最优动作。
\end{proposition}

\begin{proof}
对任意动作$\pi$，其在信念$b$下的期望损失为
\[
R(\pi,b)=\sum_{y\in\mathcal H}b(y)L(\pi,y).
\]
这是$b$的线性函数，也是有限集合$\mathcal A$上的实值函数。由$\pi^\star$的定义，
\[
R(\pi^\star,b)\le R(\pi,b),\qquad \forall \pi\in\mathcal A.
\]
因此不存在在$b$下期望损失更低的动作。若存在另一个动作$\pi'$满足$R(\pi',b)=R(\pi^\star,b)$，则$\pi'$与$\pi^\star$并列为贝叶斯最优；本文按固定顺序打破平局，不改变最优性。回答动作内部还要在三个结果桶中选择一个标签，但该选择同样是在有限集合上取最小期望损失，因而可以与动作选择合并为同一个贝叶斯决策。
\end{proof}

\section{路由效用的分解}

\begin{proposition}
设确定查询集合为$\mathcal D$，前瞻性查询集合为$\mathcal V$，路由器预测为$\hat{\mathcal D}$和$\hat{\mathcal V}$。混合效用可分解为
\[
U=\frac{1}{n}\sum_{q\in\mathcal D\cap\hat{\mathcal D}}u_D(q)
+\frac{1}{n}\sum_{q\in\mathcal V\cap\hat{\mathcal V}}u_V(q),
\]
其中$u_D(q)$为确定查询的F1，$u_V(q)=1-L(q)/10$为前瞻性查询的效用。
\end{proposition}

\begin{proof}
设$n=|\mathcal D|+|\mathcal V|$。对每条查询定义指示变量
\[
\mathbb I_D(q)=\mathbf 1[q\in\mathcal D],\qquad
\mathbb I_V(q)=\mathbf 1[q\in\mathcal V].
\]
路由器正确时获得对应分支效用，错误时获得0，因此
\[
U=\frac{1}{n}\sum_{q\in\mathcal D}\mathbb I[q\in\hat{\mathcal D}]u_D(q)
+\frac{1}{n}\sum_{q\in\mathcal V}\mathbb I[q\in\hat{\mathcal V}]u_V(q).
\]
把指示变量写成集合交集，即得
\[
U=\frac{1}{n}\sum_{q\in\mathcal D\cap\hat{\mathcal D}}u_D(q)
+\frac{1}{n}\sum_{q\in\mathcal V\cap\hat{\mathcal V}}u_V(q).
\]
oracle路由器对每条查询都正确，故
\[
U_{\text{oracle}}=\frac{1}{n}\sum_{q\in\mathcal D}u_D(q)+\frac{1}{n}\sum_{q\in\mathcal V}u_V(q).
\]
两者之差为错路由项，等于被错路由查询在本分支上的效用之和除以$n$。这一分解说明混合效用可以拆成“路由正确时的分支效用”和“错路由损失”两部分，从而把路由错误与分支错误分开诊断。
\end{proof}

\section{软锚点更新的归一化与单调性}

\begin{proposition}
设$\mathcal H$为有限结果空间，软锚点的隐含桶为$j$，似然为
\[
L_e(k)=
\begin{cases}
p_e, & k=j,\\
(1-p_e)/(|\mathcal H|-1), & k\ne j,
\end{cases}
\]
其中$p_e\in(1/|\mathcal H|,1)$。对任意先验$b(k)>0$和可靠性$\rho\in[0,1]$，更新
\[
b'(k)=\frac{b(k)L_e(k)^\rho}{\sum_{h\in\mathcal H}b(h)L_e(h)^\rho}
\]
是概率分布；当$\rho=0$时$b'=b$，当$\rho=1$时$b'$等于先验与完整似然的归一化乘积；$b'(j)/b'(k)$随$\rho$单调增大。
\end{proposition}

\begin{proof}
分母为正，因为$b(k)>0$且$L_e(k)>0$，所以$b'$非负且和为1，是概率分布。$\rho=0$时$L_e(k)^0=1$，故$b'=b$。对$k\ne j$，
\[
\frac{b'(j)}{b'(k)}
=\frac{b(j)}{b(k)}
\left(\frac{p_e}{(1-p_e)/(|\mathcal H|-1)}\right)^\rho.
\]
括号内的比值大于1（因为$p_e>1/|\mathcal H|$），所以该比值随$\rho$单调增大。可靠性越高，更新越偏向隐含桶；可靠性为0时锚点不改变信念。硬锚点可以看作$\rho\to\infty$并将非可行桶的似然置零的极限，本文用点质量加小底实现，以避免后续软锚点无法恢复。
\end{proof}

\section{可信集的最小基数}

\begin{proposition}
设$b$为$\mathcal H$上的概率分布，$R_\alpha$为满足$\sum_{k\in R}b(k)\ge 1-\alpha$的最小基数子集。把$b$按降序排列为$b_{(1)}\ge\cdots\ge b_{(|\mathcal H|)}$，则$R_\alpha$由前$m$个最大元素组成，其中$m$是使部分和达到$1-\alpha$的最小整数。
\end{proposition}

\begin{proof}
对任意大小为$m$的子集$R$，其概率质量不超过前$m$个最大元素之和。若前$m-1$个元素之和小于$1-\alpha$，则任何大小为$m-1$的子集都无法达到覆盖率，故最小基数至少为$m$。前$m$个元素之和达到$1-\alpha$，因此它们构成一个可行集，且基数等于下界$m$，所以是最小基数可信集。若存在并列的阈值元素，任选其中一个不改变覆盖率；本文按桶的固定顺序打破平局。
\end{proof}

\chapter{LLM提示词}
\label{app:prompts}

本章给出正文三个LLM调用的完整提示词。提示词只使用截止时点之前的输入，不包含已实现结果；JSON字段与正文的结构化输出一致。

\section{路由器提示词}

\begin{lstlisting}[basicstyle=\ttfamily\scriptsize,breaklines=true,columns=fullflexible,frame=single,numbers=none]
You decide whether a financial input is determinate or vague,
then route it to one of two specialists.

DETERMINATE (RAG): the input asks for a fact, number, passage,
estimate, or projection that a filing states or lets you compute
at the cutoff. A question about a future value is determinate when
the filing itself states that value or projection.

VAGUE (ANCHOR): the input asserts or asks about a future outcome
that is not realized at the cutoff. Examples include whether
reported revenue will fall below, within, or above a guidance
range, whether a product will launch, or whether a target will be
met. A management guidance statement is a forward-looking claim,
not a realized fact, even when it contains a number.

Return JSON only:
{"vague":true|false,"route":"rag|anchor",
 "confidence":0.0,"reason":"..."}

Examples:
Input: What was Entergy's net revenue change during 2015?
 -> {"vague":false,"route":"rag"}
Input: Q1 2024 revenue guidance is $12.2-13.2 billion.
 -> {"vague":true,"route":"anchor"}
Input: What did the 10-K report for segment operating profit?
 -> {"vague":false,"route":"rag"}
Input: Will Q3 revenue land within management's guidance range?
 -> {"vague":true,"route":"anchor"}
Input: Management expects second-half revenue of $1.5-1.6 billion.
 -> {"vague":true,"route":"anchor"}

Input: {query}
\end{lstlisting}

\section{锚点抽取提示词}

\begin{lstlisting}[basicstyle=\ttfamily\scriptsize,breaklines=true,columns=fullflexible,frame=single,numbers=none]
You are an AnchorEvolve adjudicator for a vague forward-looking
financial claim. Use ONLY information available at the cutoff.
Do not use the realized outcome.

The claim asks whether reported quarterly revenue will land below,
within, or above the guidance range. The guidance range defines the
three buckets; it is NOT by itself evidence that the outcome will be
within. Use the current-quarter value relative to the guidance
midpoint, the guidance width, and the observed base rates:
- mid_to_current_ratio > 1 means the midpoint requires sequential
  growth; the larger it is, the more likely below.
- mid_to_current_ratio < 1 means the midpoint implies a sequential
  decline; the smaller it is, the more likely above.
- A narrow guidance width makes within more likely; a very wide
  range makes outside buckets more likely.
- Base rates are priors, not destiny; update them with these features.

Extract explicit anchors and return JSON only:
{
  "anchors": [
    {"level":"L2",
     "type":"current_quarter_growth|guidance_width|sector_history|company_history|other",
     "relation":"bound|precedence|lead_time|duration|conservation|rate|unit_conversion|scope",
     "value":"...", "scope":"company|sector|market",
     "reliability":0.0, "hard":false,
     "implied_bucket":"below|within|above|null",
     "rationale":"..."}
  ],
  "bucket_scores":{"below":0.0,"within":0.0,"above":0.0},
  "rationale":"..."
}

Rules: a hard anchor is an explicit constraint that can rule out a
bucket (e.g. a reliable current-quarter value far above the guidance
range); most management guidance is a soft signal. reliability is in
[0,1]. bucket_scores must sum to 1.

Input:
{claim, company, ticker, sector, cutoff, target_year,
 target_quarter, guidance_low_musd, guidance_high_musd,
 guidance_mid_musd, guidance_width_pct, current_q_musd,
 current_q_reliability, observed_bucket_prior, sector_bucket_prior,
 evidence}
\end{lstlisting}

\section{阅读器提示词}

\begin{lstlisting}[basicstyle=\ttfamily\scriptsize,breaklines=true,columns=fullflexible,frame=single,numbers=none]
System:
You are a financial analyst answering a question about a long SEC
filing. Use only the excerpts below. The question may require
arithmetic; do the arithmetic carefully. If the excerpts do not
contain enough information to answer, return "insufficient" and set
sufficient to false.

Return only a JSON object with these keys:
- answer: the final answer as a short string (number or phrase)
- confidence: a number between 0 and 1
- sufficient: true if the excerpts are enough to answer, otherwise false
- missing: a short description of the missing information
- evidence_ids: the paragraph ids you used
- rationale: at most 40 words, no markdown

User:
Use only the excerpts below. Do not use outside knowledge.
This is acquisition round {round}; the excerpts include additional
paragraphs retrieved after you reported missing information.
Question: {question}
Document excerpts:
{evidence}
Return the JSON object now.
\end{lstlisting}

\section{提示词设计说明}

三个提示词共同遵循三条约束。第一，时点约束：路由器只接收查询与截止时点元数据，锚点提取器只接收截止时点之前的证据，阅读器只接收已获取的段落块。第二，结构化输出：三个调用都返回固定字段的JSON，解析失败在评价中计为错误，不进行静默修复。第三，可审计性：锚点提示词要求逐条给出类型、关系、作用域、可靠性、硬软标记、隐含桶和理由；阅读器要求给出证据编号和缺失信息，使答案可以追溯到具体段落或表格块。

\chapter{核心代码}
\label{app:code}

本章给出与正文算法直接对应的核心代码片段，完整实现和运行脚本随论文发布。代码只依赖公开的Python包和本地BM25索引，不包含API密钥。

\section{义务项与覆盖}

\begin{lstlisting}[language=Python]
import re

OBLIGATION_RE = re.compile(
    r"\$?\d[\d,]*(?:\.\d+)?%?|"
    r"\b(?:January|February|March|April|May|June|July|"
    r"August|September|October|November|December)\b",
    flags=re.I,
)

def obligation_terms(text):
    terms, seen = [], set()
    for match in OBLIGATION_RE.findall(str(text or "")):
        term = match.strip().lower().replace("$", "").replace(",", "").replace("%", "")
        if term in {"", "0", "1", "2", "3", "4", "5", "6", "7", "8", "9"}:
            continue
        if term not in seen:
            seen.add(term)
            terms.append(term)
    return terms

def uncovered_terms(terms, chunks):
    text = " ".join(str(getattr(c, "text", "")) for c in chunks)
    lowered = text.lower().replace(",", "")
    covered = set()
    for term in terms:
        if term.isdigit() and len(term) == 4 and term.startswith(("19", "20")):
            if re.search(r"(?<!\d)" + re.escape(term) + r"(?!\d)", lowered):
                covered.add(term)
        elif re.search(r"(?<![0-9.])" + re.escape(term) + r"(?![0-9])", lowered):
            covered.add(term)
    return [t for t in terms if t not in covered]
\end{lstlisting}

\section{锚点信念更新与裁决门}

\begin{lstlisting}[language=Python]
import numpy as np

BUCKETS = ["below", "within", "above"]
HARD_FLOOR = 1e-3
SOFT_P_E = 0.7
LOSS = {"over": 12.0, "under": 8.0, "wrong": 10.0,
        "miss": 8.0, "width": 2.0, "abstain": 5.0}

def normalize(v):
    v = np.maximum(np.asarray(v, dtype=float), 1e-12)
    return v / v.sum()

def belief_from_anchors(prior, parsed):
    b, anchors = prior.copy(), []
    if not parsed:
        return b, anchors
    for anchor in parsed.get("anchors", []) or []:
        implied = str(anchor.get("implied_bucket", "")).lower()
        if implied not in BUCKETS:
            continue
        rel = max(0.0, min(1.0, float(anchor.get("reliability", 0.5) or 0.5)))
        target = BUCKETS.index(implied)
        hard = bool(anchor.get("hard", False)) and rel >= 0.65
        anchors.append({"type": anchor.get("type"),
                        "relation": anchor.get("relation"),
                        "implied_bucket": implied,
                        "reliability": rel, "hard": hard,
                        "scope": anchor.get("scope")})
        if hard:
            b = np.array([1.0 if i == target else HARD_FLOOR for i in range(3)])
        else:
            like = np.full(3, (1.0 - SOFT_P_E) / 2.0)
            like[target] = SOFT_P_E
            b = b * np.power(like, rel)
        b = normalize(b)
    scores = parsed.get("bucket_scores")
    if isinstance(scores, dict):
        marginal = normalize(np.array([float(scores.get(x, 0.0)) for x in BUCKETS]))
        b = normalize(0.5 * b + 0.5 * marginal)
    return b, anchors

def credible_set(b, alpha=0.9):
    order = np.argsort(-b)
    out, total = [], 0.0
    for idx in order:
        out.append(int(idx))
        total += float(b[int(idx)])
        if total >= alpha:
            break
    return sorted(out)

def gate(b, alpha=0.9):
    exp_ans = [sum(float(b[y]) * (0.0 if a == y else
               LOSS["over"] if (a, y) == (2, 0) else
               LOSS["under"] if (a, y) == (0, 2) else
               LOSS["wrong"]) for y in range(3)) for a in range(3)]
    a = int(np.argmin(exp_ans))
    R = credible_set(b, alpha)
    exp_set = sum(float(b[y]) * ((0.0 if y in R else LOSS["miss"]) +
                 LOSS["width"] * max(0, len(R) - 1)) for y in range(3))
    if exp_ans[a] <= exp_set and exp_ans[a] <= LOSS["abstain"]:
        return {"mode": "answer", "a": a, "R": R}
    if exp_set <= LOSS["abstain"]:
        return {"mode": "adjudicate", "a": a, "R": R}
    return {"mode": "abstain", "a": a, "R": R}
\end{lstlisting}

\section{校准随机森林}

\begin{lstlisting}[language=Python]
import numpy as np
from sklearn.ensemble import RandomForestClassifier

for fold in range(FOLDS):
    train_rows = [row for row in rows
                  if fold_of(company_key(row[2]), seed) != fold
                  and row[1] < boundary]
    fold_test = [row for row in test_rows
                 if fold_of(company_key(row[2]), seed) == fold]
    X_train, C_train, y_train, meta_train = build(train_rows, by_company)
    X_test, C_test, y_test, meta_test = build(fold_test, by_company)
    train_features = np.hstack([X_train, C_train])
    test_features = np.hstack([X_test, C_test])
    model = make_rf(seed, class_weight, min_samples_leaf)
    model.fit(train_features, list(map(str, y_train)))
    probabilities = align(model.predict_proba(test_features), model.classes_)
    groups = np.asarray([company_key(record) for record in meta_train])
    oof, truths = [], []
    labels = np.asarray(list(map(str, y_train)))
    for fit_index, cal_index in GroupKFold(n_splits=min(3, len(set(groups)))).split(
        train_features, labels, groups
    ):
        inner = make_rf(seed, class_weight, min_samples_leaf)
        inner.fit(train_features[fit_index], labels[fit_index].tolist())
        oof.append(align(inner.predict_proba(train_features[cal_index]),
                         inner.classes_))
        truths.extend(labels[cal_index].tolist())
    temperature = fit_temperature(np.vstack(oof), truths)
\end{lstlisting}

\chapter{实现结构与复现边界}
\label{app:implementation}

\section{实现结构}

实现分为五个模块。路由模块读取原始查询和截止时点元数据，返回路由标签、置信度和简短理由。锚点模块读取断言、指引区间、当前季度值、基率和证据项，返回L2锚点及其可靠性、硬软标记和隐含桶。义务检索分支模块建立时点段落与表格索引，按义务覆盖迭代取证，并调用阅读器返回闭合答案。校准模块在开发集与验证集上训练随机森林并做温度缩放。汇总模块把逐条预测、路由决策和指标写入统一结果文件。各模块通过统一的数据结构连接，LLM调用按请求负载的哈希缓存，因此相同输入可以复现相同输出。

\section{复现边界}

公开结果包含逐条预测、路由决策和汇总指标；LLM响应缓存不随论文提交。随机森林的折种子、校准方式和融合权重选择在正文中给出。所有API密钥从环境变量读取，不写入结果文件或论文源文件。具名基线为同基座适配，代码、提示词和逐条输出随论文发布；原始系统在其原始数据上的未修改端到端结果不在本文的复现范围内。

% ==================== 附录 G（修订版补充结果） ====================
\chapter{全动作读出与演化版补充结果}
\label{app:evolution}

\begin{revbox}
\noindent 本附录给出修订版的全动作读出、公平基线和 DeepSeek 演化版补充结果。附录 A 的十一项指标和原表 6.2/6.3 是旧受限 gate 口径；本节表格使用全动作 gate。所有表格的测试集为 227 条 XBRL 核验断言（2022--2025），配对区间为 10,000 次 claim 级 bootstrap。
\end{revbox}

\section{全动作读出与受限 gate}

\begin{table}[H]
\centering
\zihao{5}
\setlength{\tabcolsep}{3pt}
\caption{受限 gate 与全动作 gate 的决策损失（227 条）。}
\label{tab:g-fullaction}
\begin{tabular}{lrrrr}
\toprule
基座 & Fusion 受限 & Fusion 全动作 & RF 受限 & RF 全动作 \\
\midrule
DeepSeek & 3.604 & \cellcolor{yellow!25}2.863 & 3.454 & \cellcolor{yellow!25}2.881 \\
GLM 5.3 & 3.648 & \cellcolor{yellow!25}2.956 & 3.460 & \cellcolor{yellow!25}2.885 \\
Qwen3-8B & 3.577 & \cellcolor{yellow!25}2.837 & 3.454 & \cellcolor{yellow!25}2.881 \\
\bottomrule
\end{tabular}
\end{table}

\begin{revbox}
\noindent Fusion 的全动作减损：DeepSeek $-0.740$ $[-1.031,-0.423]$；GLM $-0.690$ $[-1.000,-0.363]$；Qwen $-0.740$ $[-1.048,-0.423]$。RF 的全动作减损：DeepSeek $-0.573$ $[-0.855,-0.282]$；GLM $-0.575$ $[-0.858,-0.274]$；Qwen $-0.573$ $[-0.855,-0.282]$。六个区间均排除零。
\end{revbox}

\section{公平 gated 基线}

\begin{table}[H]
\centering
\zihao{5}
\setlength{\tabcolsep}{3pt}
\caption{Fusion 全动作 vs 最好 gated 具名基线（同一全动作 gate）。}
\label{tab:g-fair}
\begin{tabular}{llrrr}
\toprule
基座 & 最好 gated baseline & 基线损失 & Fusion 损失 & 差 [95\% CI] \\
\midrule
DeepSeek & Self-RAG-style & 3.048 & 2.863 & $-0.185$ $[-0.581,0.211]$ \\
GLM 5.3 & FinDVer-style & 3.221 & 2.956 & $-0.265$ $[-0.699,0.177]$ \\
Qwen3-8B & Fin-Fact-style & 4.846 & 2.837 & $-2.009$ $[-2.687,-1.330]$ \\
\bottomrule
\end{tabular}
\end{table}

\begin{revbox}
\noindent 公平 gate 下 Fusion 的点估计低于全部基线；Qwen 显著，DeepSeek/GLM 相对最强基线不显著。原稿的“显著低于所有基线”只成立于“baseline 只能点答案、FinPlanRAG 可以用集合/弃权”的不对等设定。
\end{revbox}

\section{DeepSeek 演化版与消融}

\begin{table}[H]
\centering
\zihao{5}
\setlength{\tabcolsep}{2.5pt}
\caption{DeepSeek 演化版系统与消融（227 条；全动作 Loss）。}
\label{tab:g-evolution}
\begin{tabular}{lrrrr}
\toprule
系统 & Acc & Bal & Macro & Loss \\
\midrule
旧论文 Fusion（受限 gate） & .524 & .471 & .468 & 3.604 \\
\rowcolor{yellow!18} 新锚点 + 演化 + 演化先验 RF（$w=.20$） & \textbf{.551} & \textbf{.489} & \textbf{.503} & \textbf{2.890} \\
\rowcolor{yellow!18} 演化先验 RF 单独（诊断） & \textbf{.555} & \textbf{.503} & \textbf{.508} & \textbf{2.828} \\
\rowcolor{yellow!18} 新锚点 + 演化 + baseline RF（$w=.05$） & .542 & .499 & .497 & 2.855 \\
\rowcolor{yellow!18} 新锚点，静态（$w=.05$） & .546 & .505 & .503 & 2.908 \\
\rowcolor{yellow!18} 旧锚点 + 演化（$w=.30$） & .520 & .439 & .438 & 2.793 \\
\rowcolor{yellow!18} 旧锚点，静态（$w=.20$） & .507 & .450 & .450 & 2.899 \\
\bottomrule
\end{tabular}
\end{table}

\begin{revbox}
\noindent 最终系统相对旧系统损失差 $-0.714$ $[-1.057,-0.352]$；相对演化先验 RF 诊断 $+0.062$ $[-0.141,0.273]$。准确率 +2.7 点，McNemar $p\approx0.69$；点估计，不作为显著优势。
\end{revbox}

\section{演化先验特征}

\begin{table}[H]
\centering
\zihao{5}
\setlength{\tabcolsep}{3pt}
\caption{演化先验特征（无 LLM；validation 185 条，test 227 条；3 seeds/300 trees）。}
\label{tab:g-prior}
\begin{tabular}{llrrrr}
\toprule
Split & 特征集 & Acc & Bal & Macro & Loss \\
\midrule
Validation & baseline & .557 & .556 & .536 & 3.784 \\
\rowcolor{yellow!18} Validation & all evolved, no artifact & \textbf{.605} & \textbf{.572} & \textbf{.566} & \textbf{3.740} \\
Test & baseline RF（受限 / 全动作） & .520 & .476 & .473 & 3.463 / 2.908 \\
\rowcolor{yellow!18} Test & evolved-prior RF（受限 / 全动作） & \textbf{.555} & \textbf{.503} & \textbf{.508} & \textbf{3.278 / 2.828} \\
\bottomrule
\end{tabular}
\end{table}

\section{Router 集成}

\begin{table}[H]
\centering
\zihao{5}
\caption{固定 router 预测、替换前瞻分支损失后的混合效用（527 条已验证子集）。}
\label{tab:g-router}
\begin{tabular}{lrr}
\toprule
分支 & 旧前瞻损失 & 新前瞻损失 \\
\midrule
Router & .389 & \cellcolor{yellow!25}.437 \\
Always-Anchor & .258 & \cellcolor{yellow!25}.306 \\
Always-RAG & .131 & .131 \\
\bottomrule
\end{tabular}
\end{table}

\section{修订版提示词与配置摘要}

\begin{revbox}
\noindent 修订版锚点提示词要求：锚点引用 \texttt{P\#}/\texttt{E\#} 来源与不晚于 cutoff 的日期；\texttt{reliability} 为隐含桶正确的校准概率；指引区间本身不作为“区间内”的证据；证据不足时返回 \texttt{missing}；公司名默认脱敏。配置项包括：硬锚点阈值 .65、软锚点 $p_e=.7$、混合权重 $\gamma=.5$、收缩系数 $\lambda=.5$、Beta 先验 $(1,1)$、损失 $(12,8,10,8,2,5)$、融合权重网格 $\{0,.05,.1,.2,.3,.5,1\}$、RF 600 棵树/叶节点 5/5 seeds。完整 prompt、配置和逐条结果随代码发布。
\end{revbox}

\begin{revbox}
\noindent\textbf{旧附录说明。}附录 A 的十一项指标、原表 6.2/6.3 以及旧锚点消融是\textbf{旧受限 gate} 口径，保留用于历史对照。修订版论文的主结果以本节和表\ref{tab:fullaction}、表\ref{tab:evolution}、表\ref{tab:prior} 为准。
\end{revbox}


% ==================== 致谢（终稿恢复） ====================
% 按学院通知：预答辩与查重阶段不排入致谢，终稿阶段取消注释并补写致谢内容。
% \clearpage
% \sufehead{致\quad 谢}

% ==================== 个人简历（终稿恢复） ====================
% 按学院通知：预答辩与查重阶段不排入个人简历及在学期间发表的研究成果，终稿阶段取消注释。
% \clearpage
% \sufehead{个人简历及在学期间发表的研究成果}
% {\zihao{5}\songti
% \setlength{\baselineskip}{16pt}
% \noindent 姓名：高深深\\
% 出生年月：\underline{\hspace{2.4cm}}年\underline{\hspace{1.2cm}}月\underline{\hspace{1.2cm}}日\\
% 本科毕业院校：\underline{\hspace{6cm}}\\
% 本科专业：\underline{\hspace{6cm}}\\
% 获学士学位时间：\underline{\hspace{2.4cm}}年\underline{\hspace{1.2cm}}月\\[0.4em]
% 硕士研究生：上海财经大学金融学院，金融工程与量化投资专业，2025年9月入学，导师罗丹。\\[0.6em]
% \textbf{在学期间发表的研究成果：}无。
% }

% ==================== 封底 ====================
\clearpage
\thispagestyle{empty}
\vspace*{\fill}
\begin{center}
{\zihao{3}\heiti 上海财经大学}\\[1.2em]
{\zihao{4}\songti 上海市国定路 777 号}
\end{center}
\vspace*{\fill}


\end{document}
